\documentclass[11pt]{article}
\usepackage[T1]{fontenc}
\usepackage[utf8]{inputenc}
\usepackage{lmodern}
\usepackage[a4paper,margin=27mm]{geometry}
\usepackage{amsmath,amssymb,amsthm,mathtools,aliascnt}
\usepackage{microtype,booktabs,longtable,array}
\usepackage{xcolor}
\usepackage{enumitem}
\usepackage{listings}
\usepackage{hyperref}
\hypersetup{colorlinks=true,linkcolor=blue!45!black,citecolor=blue!45!black,urlcolor=blue!45!black,
 pdftitle={The Baker--Campbell--Hausdorff Formula: Complete Proofs and All-Order Expansions},
 pdfauthor={}}
\usepackage[nameinlink,noabbrev]{cleveref}
\setlength{\emergencystretch}{2em}
\setlength{\parskip}{3pt}
\setcounter{tocdepth}{2}
\numberwithin{equation}{section}
\newtheorem{theorem}{Theorem}[section]
\newaliascnt{lemma}{theorem}
\newtheorem{lemma}[lemma]{Lemma}
\aliascntresetthe{lemma}
\newaliascnt{proposition}{theorem}
\newtheorem{proposition}[proposition]{Proposition}
\aliascntresetthe{proposition}
\newaliascnt{corollary}{theorem}
\newtheorem{corollary}[corollary]{Corollary}
\aliascntresetthe{corollary}
\theoremstyle{definition}
\newaliascnt{definition}{theorem}
\newtheorem{definition}[definition]{Definition}
\aliascntresetthe{definition}
\newaliascnt{example}{theorem}
\newtheorem{example}[example]{Example}
\aliascntresetthe{example}
\theoremstyle{remark}
\newaliascnt{remark}{theorem}
\newtheorem{remark}[remark]{Remark}
\aliascntresetthe{remark}
\DeclareMathOperator{\ad}{ad}
\DeclareMathOperator{\Ad}{Ad}
\DeclareMathOperator{\tr}{tr}
\DeclareMathOperator{\BCH}{BCH}
\DeclareMathOperator{\GL}{GL}
\DeclareMathOperator{\End}{End}
\DeclareMathOperator{\Span}{span}
\DeclareMathOperator{\spec}{spec}
\DeclareMathOperator{\id}{id}
\DeclareMathOperator{\diag}{diag}
\newcommand{\kk}{\mathbb{k}}
\newcommand{\R}{\mathbb{R}}
\newcommand{\C}{\mathbb{C}}
\newcommand{\Q}{\mathbb{Q}}
\newcommand{\N}{\mathbb{N}}
\newcommand{\Z}{\mathbb{Z}}
\newcommand{\g}{\mathfrak{g}}
\newcommand{\h}{\mathfrak{h}}
\newcommand{\cL}{\mathcal{L}}
\newcommand{\cA}{\mathcal{A}}
\newcommand{\cT}{\mathcal{T}}
\newcommand{\eps}{\varepsilon}
\newcommand{\norm}[1]{\left\lVert#1\right\rVert}
\newcommand{\coeff}[2]{[#1]#2}
\newcommand{\RB}[1]{[#1]_{\mathrm R}}
\newcommand{\ord}{\operatorname{ord}}
\newcommand{\dt}{\,dt}
\newcommand{\dd}{\,d}
\newcommand{\one}{\mathbf{1}}
\lstset{basicstyle=\small\ttfamily,columns=fullflexible,breaklines=true,frame=single,
 showstringspaces=false,commentstyle=\itshape\color{black!60}}
\title{The Baker--Campbell--Hausdorff Formula:\\
 Complete Proofs, All-Order Expansions,\\and Analytic Qualifications}
\author{}
\date{16 September 2026}
\begin{document}
\maketitle
\begin{abstract}
We give a systematic proof of the mathematical assertions and displayed
identities in the Wikipedia article on the Baker--Campbell--Hausdorff formula,
with necessary corrections to its geometric and operator-theoretic
qualifications. The treatment separates formal identities, convergent local
identities, and global exponential identities. A direct proof of the
Dynkin--Specht--Wever projection proves Friedrichs' characterization of
primitive series and yields Dynkin's formula with every coefficient specified.
An independent differential proof produces the Bernoulli recursion and the
Poincar\'e integral formula. Explicit finite word sums describe every
homogeneous and bihomogeneous BCH component. For the Zassenhaus product we
prove existence, uniqueness, two recursive constructions, a nonrecursive
weighted-tree formula, and a quantitative convergence criterion. We also
prove the local Lie correspondence, product formulas of Trotter and Suzuki,
all-order Maurer--Cartan and metric expansions, and the quantum-mechanical
identities in concrete Hilbert-space realizations. Counterexamples clarify
logarithm branches, nonconvergence, degenerate pullbacks, and the distinction
between infinitesimal and exponentiated commutation relations. Exact rational
computations through degree ten accompany the proofs.
\end{abstract}
\begingroup
\small
\setlength{\parskip}{0pt}
\tableofcontents
\endgroup
\newpage

\section{Scope, conventions, and the three meanings of BCH}\label{sec:scope}
The central problem is to describe a logarithm of a product of exponentials:
\begin{equation}\label{eq:main}
 e^X e^Y=e^{Z(X,Y)}.
\end{equation}
The reference page is the revision of the Wikipedia article
\emph{Baker--Campbell--Hausdorff formula} with revision identifier
\texttt{1368909113}, last edited on 11 August 2026 and consulted on
16 September 2026 \cite{wiki}. The coverage table in
\cref{sec:coverage} identifies where each family of assertions is proved.
Historical attributions and lists of links are not mathematical theorems;
likewise, a theorem merely named as an application is not thereby stated on
the page. In particular, the Stone--von Neumann classification theorem and
the Golden--Thompson inequality are not additional claims that the present
paper purports to prove.

There are three distinct settings.

\paragraph{Formal algebra.}
$X$ and $Y$ have positive degree. Every identity is an identity of formal
series, and every fixed-degree coefficient is a finite sum. No numerical
convergence is intended.
\paragraph{Local analysis.}
$X,Y$ belong to a Banach algebra or a finite-dimensional Lie algebra and
are sufficiently small. The series converges and specifies the logarithm
branch issuing from zero.
\paragraph{Global identities under special relations.}
Commutation relations may allow the series to terminate or to be replaced
by a closed expression. An independent differential argument can then prove
an exponential identity beyond the convergence domain of its Taylor series.

Conflating these settings is the source of several apparent paradoxes.
A matrix may have logarithms even though its BCH series diverges; a closed
expression may continue the local BCH logarithm without making the original
series convergent; a logarithm need not belong to the Lie algebra generated
by its two inputs.

Throughout, $[A,B]=AB-BA$ in an associative algebra and
$\ad_A B=[A,B]$. Composition of adjoint maps is read from right to left.
For a nonempty word $w=a_1\cdots a_n$ define its \emph{right-nested bracket}
by
\begin{equation}\label{eq:rightbracket}
 \RB{a_1\cdots a_n}=[a_1,[a_2,\ldots,[a_{n-1},a_n]\ldots]],
 \qquad \RB{a}=a.
\end{equation}
Thus $\RB{XXY}=[X,[X,Y]]$; a superscript inside a word denotes repeated
letters, not an associative power after the bracketing operation has been
performed. We write $Z_n$ for total degree $n$, and $Z_{p,q}$ for degree $p$
in $X$ and degree $q$ in $Y$.

The scalar functions that organize the entire subject are
\begin{equation}\label{eq:scalarfunctions}
 \phi_-(z)=\frac{1-e^{-z}}{z},\qquad
 \phi_+(z)=\frac{e^z-1}{z},\qquad
 f(z)=\frac{z}{1-e^{-z}},\qquad
 \psi(x)=\frac{x\log x}{x-1}.
\end{equation}
The first two are entire, with value $1$ at zero. The function $f$ is
meromorphic with a removable singularity at zero and poles at the nonzero
points $2\pi i\Z$. The function $\psi$ is used near $x=1$, with a logarithm
branch analytic there. Expressions such as $\phi_-(M)$ always denote
functional calculus, not an assumption that $M$ is invertible.

\section{Formal series and a coefficient formula for every word}\label{sec:words}
\subsection{The completed associative algebra}
Let $\kk$ be a field of characteristic zero and $V=\kk X\oplus\kk Y$.
The tensor algebra and its degree completion are
\[
 T(V)=\bigoplus_{n\geq0}V^{\otimes n},\qquad
 \widehat T(V)=\prod_{n\geq0}V^{\otimes n}.
\]
Multiplication is concatenation, extended by the Cauchy product. The
filtration $F^m\widehat T=\prod_{n\geq m}V^{\otimes n}$ satisfies
$F^mF^n\subseteq F^{m+n}$. Consequently, substitution of positive-degree
series into another formal series is well defined.

\begin{lemma}[Formal exponential and logarithm]\label{lem:explog}
The maps
\begin{equation}\label{eq:formalexplog}
 \exp A=\sum_{n\geq0}\frac{A^n}{n!},\qquad
 \log(1+U)=\sum_{k\geq1}\frac{(-1)^{k-1}}kU^k
\end{equation}
are mutually inverse bijections between $F^1\widehat T$ and
$1+F^1\widehat T$. If $A$ and $B$ commute, then
$e^{A+B}=e^Ae^B$.
\end{lemma}
\begin{proof}
The series are defined because only finitely many powers contribute in
any fixed degree. The scalar formal identities
$\log(e^t)=t$ and $e^{\log(1+t)}=1+t$ follow, for example, by
formal differentiation and comparison of constant terms. Substituting
$A$ or $U$ is legitimate: powers of one fixed algebra element commute
with each other, regardless of whether the ambient algebra is
commutative. If $AB=BA$, the binomial theorem gives
$(A+B)^n=\sum_{j=0}^n\binom nj A^jB^{n-j}$; summing proves the last
assertion. No commutativity of unrelated elements has been used.
\end{proof}

Define
\begin{equation}\label{eq:defbch}
 Z(X,Y)=\BCH(X,Y)=\log(e^Xe^Y)=\sum_{n\geq1}Z_n(X,Y).
\end{equation}
By \cref{lem:explog}, this is the unique positive-degree formal solution
of \eqref{eq:main}. Expanding the two exponentials first gives
\begin{equation}\label{eq:associativeblocks}
 Z(X,Y)=\sum_{k\geq1}\frac{(-1)^{k-1}}k
 \sum_{\substack{r_i,s_i\geq0\\r_i+s_i>0\ (1\leq i\leq k)}}
 \frac{X^{r_1}Y^{s_1}\cdots X^{r_k}Y^{s_k}}
 {\prod_{i=1}^k r_i!s_i!}.
\end{equation}
This proves the associative expansion on the reference page. It does not
yet prove that its homogeneous parts are Lie polynomials.

\subsection{A finite nonrecursive formula}
For a nonempty word $w$ in $X,Y$, put
\begin{equation}\label{eq:blockweight}
 a(w)=
 \begin{cases}
 (r!s!)^{-1},&w=X^rY^s,\quad r+s=|w|,\\
 0,&\text{otherwise}.
 \end{cases}
\end{equation}
Equivalently, $a(w)=0$ precisely when $w$ contains the adjacent pattern
$YX$. Pure powers are permitted.

\begin{theorem}[Every associative BCH coefficient]\label{thm:wordcuts}
For $w=a_1\cdots a_n$, the coefficient $c(w)$ of $w$ in $Z$ is
\begin{equation}\label{eq:wordcuts}
 c(w)=\sum_{k=1}^{n}\frac{(-1)^{k-1}}k
 \sum_{0=i_0<i_1<\cdots<i_k=n}
 \prod_{j=1}^{k}a(a_{i_{j-1}+1}\cdots a_{i_j}).
\end{equation}
Thus
\begin{equation}\label{eq:wordexpansion}
 Z_n=\sum_{w\in\{X,Y\}^n}c(w)w.
\end{equation}
For a prescribed word, the total number of cut patterns before zero
terms are discarded is $2^{n-1}$.
\end{theorem}
\begin{proof}
Write $U=e^Xe^Y-1$. Its nonempty-word coefficient is $a(w)$.
A contribution to the coefficient of $w$ in $U^k$ is exactly a cut of
$w$ into $k$ nonempty contiguous blocks; the contribution is the product
of their coefficients. Substituting into the logarithm proves
\eqref{eq:wordcuts}. Choosing a subset of the $n-1$ gaps is the same
thing as choosing a cut pattern, proving the count.
\end{proof}
For example, $c(XY)=1-1/2=1/2$, while $c(YX)=-1/2$.
For $XXY$, the one-block contribution is $1/2$, the two-block
contribution is $-\tfrac12(1+\tfrac12)$, and the three-block contribution
is $1/3$. Their sum is $1/12$. This small computation already explains
why the coefficients are rational and why cancellation matters.

\subsection{Any finite number of exponential factors}\label{sec:multifactor}
The same construction works without change for
$\log(e^{A_1}\cdots e^{A_m})$. In the alphabet $A_1,\ldots,A_m$, replace
$a(w)$ by $(r_1!\cdots r_m!)^{-1}$ if
$w=A_1^{r_1}\cdots A_m^{r_m}$, and by zero otherwise. If the factors are
$e^{\alpha_j X}$ or $e^{\beta_jY}$ with letters repeated in different
positions, first use a distinct letter for each factor and then substitute.
The resulting coefficients are finite polynomials in the scalar factor
weights. This observation supplies all orders, not just leading errors,
for the splitting formulas in \cref{sec:splitting}.

\section{Hopf algebra, primitive elements, and Dynkin's formula}\label{sec:hopf}
\subsection{The coproduct must take values in a completion}
The relevant target of the coproduct on infinite series is
\[
 \widehat T\,\widehat\otimes\,\widehat T
 =\prod_{n\geq0}\bigoplus_{p+q=n}V^{\otimes p}\otimes V^{\otimes q}.
\]
The ordinary algebraic tensor product is generally too small. Define the
continuous algebra homomorphism $\Delta$ by
\begin{equation}\label{eq:coproductgenerators}
 \Delta X=X\otimes1+1\otimes X,\qquad
 \Delta Y=Y\otimes1+1\otimes Y.
\end{equation}
For a word $w=a_1\cdots a_n$, let $w_I$ be the subword at the positions
$I\subseteq\{1,\ldots,n\}$, in their original order. Then
\begin{equation}\label{eq:unshuffle}
 \Delta w=\sum_{I\subseteq\{1,\ldots,n\}}w_I\otimes w_{I^c}.
\end{equation}
The empty subword is $1$. Formula \eqref{eq:unshuffle} follows by
multiplying the $n$ factors \eqref{eq:coproductgenerators}; it also proves
coassociativity, because an iterated coproduct assigns each letter to
one of three tensor factors independently of the order of assignment.

Let $\eps$ extract the constant term. Define the antipode on words by
\begin{equation}\label{eq:antipode}
 S(a_1\cdots a_n)=(-1)^n a_n\cdots a_1,\qquad S(1)=1.
\end{equation}
It is an antiautomorphism. The convolution identities are
\begin{equation}\label{eq:antipodeidentities}
 m(\id\otimes S)\Delta w=m(S\otimes\id)\Delta w=\eps(w)1.
\end{equation}
Here $m$ is multiplication. For completeness, the first identity can be
proved inductively. If $a$ is a letter and
$\Delta v=\sum v_{(1)}\otimes v_{(2)}$, then the sum for $av$ is
\[
 a\sum v_{(1)}S(v_{(2)})-
 \sum v_{(1)}S(v_{(2)})a,
\]
which vanishes by the induction hypothesis. The degree-zero case is $1$.
The second identity follows in the same way, or by applying reversal.
Thus $T(V)$, and continuously its completion, has the asserted Hopf
algebra structure.

\begin{definition}
A positive-degree series $P$ is \emph{primitive} if
$\Delta P=P\otimes1+1\otimes P$. A series $G$ with constant term $1$ is
\emph{grouplike} if $\Delta G=G\otimes G$.
\end{definition}
The normalization for a grouplike series excludes the zero series, which
would otherwise satisfy the displayed coproduct equation but would not
belong to a group.

\begin{lemma}[Exponentials of primitives]\label{lem:grouplike}
$e^P$ is grouplike if and only if $P$ is primitive. Grouplike series form
a group under multiplication.
\end{lemma}
\begin{proof}
If $P$ is primitive, the commuting elements $P\otimes1$ and $1\otimes P$
give
\[
 \Delta e^P=e^{\Delta P}=e^{P\otimes1+1\otimes P}=e^P\otimes e^P.
\]
Conversely, if $G=e^P$ is grouplike, then
$G\otimes G=(G\otimes1)(1\otimes G)$. These two factors commute, so
\[
 \Delta P=\log(\Delta G)
 =\log(G\otimes1)+\log(1\otimes G)=P\otimes1+1\otimes P.
\]
Multiplicativity of $\Delta$ proves closure under products. A series
with constant term $1$ is invertible by its geometric series, and
$\Delta(G^{-1})=(G\otimes G)^{-1}=G^{-1}\otimes G^{-1}$.
\end{proof}

\subsection{An elementary projection proof of Friedrichs' theorem}
Let $N$ be the grading operator, $N(w)=|w|w$. Define
\begin{equation}\label{eq:dynkinoperator}
 D=m(N\otimes S)\Delta.
\end{equation}
This particular order of the two convolution factors is what produces
\emph{right}-nested rather than left-nested brackets.

\begin{lemma}[Dynkin--Specht--Wever]\label{lem:dsw}
For every nonempty word $w$, $D(w)=\RB{w}$. If $P$ is a homogeneous
primitive element of degree $n\geq1$, then
\begin{equation}\label{eq:dswprimitive}
 D(P)=nP.
\end{equation}
\end{lemma}
\begin{proof}
For a letter $a$, the definition gives $D(a)=a$. Using the same
coproduct expansion as above, and writing $|v_{(1)}|$ for degree, gives
\begin{align*}
 D(av)
 &=\sum (|v_{(1)}|+1)a v_{(1)}S(v_{(2)})
   -\sum |v_{(1)}|v_{(1)}S(v_{(2)})a\\
 &=aD(v)-D(v)a+a\eps(v).
\end{align*}
For nonempty $v$, this is $[a,D(v)]$, proving the first assertion by
induction. If $P$ is primitive, only its two primitive coproduct terms
contribute, and $N(1)=0$, so
$D(P)=N(P)S(1)+N(1)S(P)=nP$.
\end{proof}

Let $\cL\subseteq T(V)$ be the Lie subalgebra generated by $X,Y$, and let
$\cL_n=\cL\cap V^{\otimes n}$.
\begin{theorem}[Friedrichs' characterization]\label{thm:friedrichs}
The primitive series are precisely $\prod_{n\geq1}\cL_n$, the
formal degreewise sums of Lie polynomials.
\end{theorem}
\begin{proof}
The generators are primitive. A commutator of primitives is primitive,
since the mixed tensor factors commute and cancel in
$[\Delta P,\Delta Q]$. Hence every Lie polynomial, and every
convergent-in-degree sum of them, is primitive.
Conversely, each homogeneous part $P_n$ of a primitive series is
primitive. By \cref{lem:dsw},
\[
 P_n=\frac1nD(P_n)=\frac1n\sum_{|w|=n}[w]P_n\,\RB{w}\in\cL_n.
\]
Characteristic zero is essential for division by $n$.
\end{proof}
Here and below $[w]P$ means the associative coefficient of the word $w$
in $P$, not a commutator.

\begin{corollary}[Formal BCH theorem]\label{cor:formalBCH}
Every $Z_n(X,Y)$ is a homogeneous Lie polynomial. In particular
$Z-X-Y$ belongs degreewise to the Lie ideal generated by $[X,Y]$.
\end{corollary}
\begin{proof}
$X,Y$ are primitive, their exponentials are grouplike, the product is
grouplike, and its logarithm is primitive. Apply
\cref{thm:friedrichs}. A right-nested bracket of length at least two has
innermost bracket either zero or $\pm[X,Y]$, proving the final statement.
\end{proof}

\subsection{Dynkin's explicit all-orders formula}
Applying $D/n$ to the total-degree-$n$ part of
\eqref{eq:associativeblocks} proves the following formula, traditionally
associated with Dynkin \cite{dynkin}.
\begin{theorem}[Dynkin formula]\label{thm:dynkin}
In the completed free Lie algebra,
\begin{equation}\label{eq:dynkin}
 \boxed{\displaystyle
 Z(X,Y)=\sum_{k\geq1}\frac{(-1)^{k-1}}k
 \sum_{\substack{r_i,s_i\geq0\\r_i+s_i>0\ (1\leq i\leq k)}}
 \frac{\RB{X^{r_1}Y^{s_1}\cdots X^{r_k}Y^{s_k}}}
 {\bigl(\sum_{i=1}^k(r_i+s_i)\bigr)\prod_{i=1}^k r_i!s_i!}.}
\end{equation}
In each total degree the sum is finite. A summand vanishes if
$s_k>1$, or if $s_k=0$ and $r_k>1$.
\end{theorem}
\begin{proof}
The projection argument proves the formula. The stated zero cases have
innermost bracket $[Y,Y]$ or $[X,X]$ respectively. These are sufficient
zero conditions, not a complete characterization of zero brackets.
\end{proof}
Combining this with the cut formula gives a particularly compact answer
to the request for \emph{all} terms:
\begin{equation}\label{eq:allbidegrees}
 \boxed{\displaystyle
 Z_{p,q}=\frac1{p+q}
 \sum_{\substack{w\in\{X,Y\}^{p+q}\\\#_Xw=p,\ \#_Yw=q}}
 c(w)\RB{w},\qquad p+q\geq1,}
\end{equation}
where $c(w)$ is the finite sum \eqref{eq:wordcuts}. There are no
unspecified coefficients or auxiliary recurrences in
\eqref{eq:allbidegrees}. The brackets form a spanning family rather than
a basis, so different displayed terms can represent the same Lie element.
This is redundancy, not ambiguity in $Z$.

\subsection{Free Lie algebras and universal enveloping algebras}
The Lie algebra $\cL$ is free on $X,Y$, and its universal enveloping
algebra is canonically $T(V)$. Here is a precise justification, including
the point needed to substitute into an arbitrary Lie algebra.

The Poincar\'e--Birkhoff--Witt injectivity statement, proved in
\cref{sec:pbw}, embeds any Lie algebra $\h$ into $U(\h)$. Given
$x,y\in\h$, the associative homomorphism $T(V)\to U(\h)$ determined
by $X\mapsto x,Y\mapsto y$ sends every iterated commutator into $\h$.
Its restriction therefore gives a Lie homomorphism $\cL\to\h$, unique
because $X,Y$ generate $\cL$. This is exactly the universal property
of the free Lie algebra.

The inclusion $\cL\to T(V)$ extends to $U(\cL)\to T(V)$, and the
two generators give a homomorphism in the other direction. The two
composites fix $X,Y$; both associative algebras are generated by those
elements, so the composites are identities. The Hopf structures agree
on the generators, hence everywhere. This proves the enveloping-algebra
assertion on the reference page.

For arbitrary $x,y\in\h$ the safe universal expression is
$\BCH(tx,ty)\in\h[[t]]$. An infinite sum in the uncompleted vector
space $\h$ has no meaning until a topology, a nilpotence hypothesis,
or a convergence argument has been supplied.

\section{Conjugation and differentiation of the exponential}\label{sec:dexp}
\subsection{Campbell's identity and the full conjugation series}
\begin{theorem}[Campbell--Hadamard identity]\label{thm:hadamard}
Formally, or for elements of a Banach algebra,
\begin{equation}\label{eq:hadamard}
 e^{tX}Ye^{-tX}=e^{t\ad_X}Y
 =\sum_{n\geq0}\frac{t^n}{n!}\ad_X^nY.
\end{equation}
Equivalently, $\Ad_{e^X}=e^{\ad_X}$.
\end{theorem}
\begin{proof}
Set $F(t)=e^{tX}Ye^{-tX}$. Direct differentiation gives
$F'(t)=XF(t)-F(t)X=\ad_XF(t)$ and $F(0)=Y$.
The coefficient recursion uniquely solves this equation as the displayed
series. In a Banach algebra,
$\norm{\ad_X^nY}\leq2^n\norm X^n\norm Y$, so the series converges
absolutely for every $t$ and the same differentiation is legitimate.
\end{proof}
An entirely associative version is
\begin{equation}\label{eq:adbinomial}
 \ad_X^nY=\sum_{j=0}^n(-1)^j\binom nj X^{n-j}YX^j.
\end{equation}
Indeed, left and right multiplication by $X$ commute, so their difference
can be raised to the $n$th power using the ordinary binomial theorem.
This proves every term of the Campbell expansion, including its
iterated-commutator notation $[X,Y]_n$ on the reference page.

Conjugation is an algebra automorphism and therefore commutes with
exponentiation. It follows that
\begin{align}
 e^Xe^Ye^{-X}&=\exp(e^XYe^{-X})
              =\exp(e^{\ad_X}Y),\label{eq:conjexp}\\
 e^Xe^Y&=\exp\left(\sum_{n\geq0}\frac{\ad_X^nY}{n!}\right)e^X.
 \label{eq:braiding}
\end{align}
These are the general braiding identities. Unlike a general BCH series,
the adjoint exponential in \eqref{eq:braiding} is an entire series for
bounded operators.

\subsection{The differential and its inverse}
\begin{theorem}[Differential of the exponential]\label{thm:dexp}
For an increment $H$,
\begin{align}
 (d\exp)_X(H)
 &=\sum_{n\geq1}\frac1{n!}\sum_{j=0}^{n-1}X^jHX^{n-1-j}
   =\int_0^1 e^{(1-s)X}He^{sX}\,ds,\label{eq:dexpintegral}\\
 e^{-X}(d\exp)_X(H)
 &=\phi_-(\ad_X)H
   =\sum_{n\geq0}\frac{(-1)^n}{(n+1)!}\ad_X^nH,\label{eq:leftdexp}\\
 (d\exp)_X(H)e^{-X}
 &=\phi_+(\ad_X)H
   =\sum_{n\geq0}\frac1{(n+1)!}\ad_X^nH.\label{eq:rightdexp}
\end{align}
All three series converge for arbitrary $X,H$ in a Banach algebra.
\end{theorem}
\begin{proof}
Differentiating $(X+uH)^n$ at $u=0$ inserts $H$ in each possible
position, giving the double sum. The coefficient of $X^pHX^q$ in the
integral is
\[
 \frac1{p!q!}\int_0^1(1-s)^p s^q\,ds=\frac1{(p+q+1)!}.
\]
The integral identity follows; the elementary beta integral here follows
by repeated integration by parts. Left multiplication by $e^{-X}$ and
\cref{thm:hadamard} give \eqref{eq:leftdexp}. Right multiplication gives
\eqref{eq:rightdexp}. The factorial denominators and the adjoint norm
bound justify termwise operations in the analytic setting.
\end{proof}
The same statements hold intrinsically for finite-dimensional Lie
groups, where multiplication by $e^{-X}$ means left or right translation
of tangent vectors; an intrinsic derivation is supplied in
\cref{sec:localgroups}.

\subsection{Bernoulli coefficients, including the convention at degree one}
Define the rational numbers $b_n$ by
\begin{equation}\label{eq:bernoulligen}
 f(z)=\frac{z}{1-e^{-z}}=\sum_{n\geq0}b_nz^n
 =\sum_{n\geq0}\frac{B_n^+}{n!}z^n.
\end{equation}
Thus $B_1^+=+1/2$, the convention used in the reference page's note.
The alternative convention for $z/(e^z-1)$ has $B_1^-=-1/2$ and
$B_n^-=(-1)^nB_n^+$.

\begin{proposition}[Complete Bernoulli coefficients]\label{prop:bernoulli}
$b_0=1$, and for $n\geq1$,
\begin{align}
 b_n&=-\sum_{j=1}^n\frac{(-1)^j}{(j+1)!}b_{n-j},\label{eq:bernrec}\\
 b_n&=\sum_{k=1}^n(-1)^{n+k}
 \sum_{\substack{j_1+\cdots+j_k=n\\j_i\geq1}}
 \prod_{i=1}^k\frac1{(j_i+1)!}.\label{eq:bernfinite}
\end{align}
Moreover,
\begin{equation}\label{eq:coth}
 f(z)=\frac z2\left(1+\coth\frac z2\right)
 =1+\frac z2+\frac{z^2}{12}-\frac{z^4}{720}
   +\frac{z^6}{30240}-\frac{z^8}{1209600}+\cdots.
\end{equation}
All odd coefficients beyond $b_1$ vanish, and the Taylor series of $f$
has radius $2\pi$.
\end{proposition}
\begin{proof}
The identity $\phi_-(z)f(z)=1$ gives \eqref{eq:bernrec} by coefficient
comparison. Expanding
$(1+\sum_{j\geq1}(-1)^jz^j/(j+1)!)^{-1}$ as a geometric series gives
\eqref{eq:bernfinite}. Algebraic simplification gives \eqref{eq:coth}.
The function $f(z)-z/2$ is even, proving the odd-coefficient assertion.
The nonzero zeros of $1-e^{-z}$ are $2\pi i\Z\setminus\{0\}$;
the closest are genuine poles at $\pm2\pi i$. The Taylor radius is
therefore exactly $2\pi$.
\end{proof}
In particular, whenever inversion is legitimate,
\begin{equation}\label{eq:inversedexp}
 H=f(\ad_X)\bigl(e^{-X}(d\exp)_X(H)\bigr).
\end{equation}
The quotient notation in this expression denotes its analytic value,
not division by an invertible $\ad_X$.

\section{Differential equations, the integral formula, and all bidegrees}
\label{sec:differentialbch}
\subsection{The differential equation for one moving factor}
\begin{theorem}\label{thm:bchode}
In the degree completion, or on a sufficiently small analytic
neighborhood of $(0,0)$, the function
$Z(t)=\BCH(X,tY)$ is characterized by
\begin{equation}\label{eq:bchode}
 Z'(t)=f(\ad_{Z(t)})Y,
 \qquad Z(0)=X.
\end{equation}
\end{theorem}
\begin{proof}
For $G(t)=e^Xe^{tY}$, its left logarithmic derivative is
$G^{-1}G'=Y$. Substitute $G=e^{Z(t)}$ in \eqref{eq:leftdexp}:
$Y=\phi_-(\ad_Z)Z'$. Its formal inverse is $f(\ad_Z)$, proving
\eqref{eq:bchode}. Conversely, that equation implies that $e^{Z(t)}$
has the same left logarithmic derivative and initial value as $G(t)$.
The equation $G'=GY$ has a unique solution, coefficient by coefficient
formally, or by the ordinary uniqueness theorem analytically.
\end{proof}
This is also a useful independent mechanism for constructing the BCH
series: its right-hand side involves only iterated Lie brackets.

\begin{theorem}[Poincar\'e integral formula]\label{thm:poincare}
With the same interpretations,
\begin{equation}\label{eq:poincare}
 \BCH(X,Y)=X+\int_0^1
 \psi\bigl(e^{\ad_X}e^{t\ad_Y}\bigr)Y\dt,
\end{equation}
where the all-orders expansion at $x=1$ is
\begin{equation}\label{eq:psiseries}
 \psi(x)=1-\sum_{n\geq1}\frac{(1-x)^n}{n(n+1)}.
\end{equation}
In particular,
\begin{equation}\label{eq:psibernoulli}
 \psi(e^z)=f(z)=\sum_{n\geq0}\frac{B_n^+}{n!}z^n.
\end{equation}
\end{theorem}
\begin{proof}
Hadamard's identity and composition of conjugations give
\[
 e^{\ad_{Z(t)}}=\Ad(e^{Z(t)})
 =\Ad(e^Xe^{tY})=e^{\ad_X}e^{t\ad_Y}.
\]
The formal logarithm of either side is $\ad_{Z(t)}$; analytically this
statement holds on the logarithm branch near the identity. Since
$\psi(e^z)=z/(1-e^{-z})$, equation \eqref{eq:bchode} becomes the
integrand of \eqref{eq:poincare}. Integrate from $0$ to $1$.
For completeness, set $u=1-x$. Multiplying
$\log(1-u)=-\sum_{n\geq1}u^n/n$ by $(1-u)/(-u)$ gives
\eqref{eq:psiseries}, initially for $|u|<1$, and also as a formal
identity. Equation \eqref{eq:psibernoulli} then follows directly.
In the formal integral, integration means integrating each polynomial
coefficient in $t$, so no interchange of improper limits is involved.
\end{proof}
There is no need for an invertible $\ad_Z$ in this argument. Likewise,
$\psi(E)$ near $E=I$ is defined by its power series in $E-I$, not by
multiplication by a nonexistent inverse of $E-I$.

\subsection{Every power of the second variable}
The notation $O(Y^2)$ hides a whole hierarchy which can be specified
explicitly. Write
\begin{equation}\label{eq:Yexpansion}
 \BCH(X,\eps Y)=X+\sum_{m\geq1}\eps^m V_m(X,Y),\qquad V_0=X.
\end{equation}
Each $V_m$ is a formal series in $X$, homogeneous of degree $m$ in $Y$.

\begin{proposition}[All powers of $Y$]\label{prop:Yall}
With $b_k=B_k^+/k!$,
\begin{equation}\label{eq:Yrecurrence}
 mV_m=\sum_{k\geq0}b_k
 \sum_{\substack{j_1+\cdots+j_k=m-1\\j_\ell\geq0}}
 \ad_{V_{j_1}}\cdots\ad_{V_{j_k}}Y.
\end{equation}
For $k=0$ the inner expression is $Y$ when $m=1$, and is zero
otherwise. In particular,
\begin{align}
 V_1&=f(\ad_X)Y
 =\frac{\ad_X}{2}\left(1+\coth\frac{\ad_X}{2}\right)Y,
 \label{eq:linearY}\\
 V_2&=\frac12\sum_{k\geq1}b_k\sum_{\ell=0}^{k-1}
 \ad_X^\ell\ad_{V_1}\ad_X^{k-1-\ell}Y.
 \label{eq:quadraticY}
\end{align}
Every $V_m$ is thereby determined by $V_0,\ldots,V_{m-1}$.
A finite, nonrecursive alternative is
\begin{equation}\label{eq:Yclosed}
 V_m=\sum_{p\geq0}\frac1{p+m}
 \sum_{\substack{w\in\{X,Y\}^{p+m}\\\#X(w)=p,\ \#Y(w)=m}}
 c(w)\RB{w},
\end{equation}
where each $c(w)$ is the finite sum \eqref{eq:wordcuts}.
\end{proposition}
\begin{proof}
Expand the right-hand side of \eqref{eq:bchode} in powers of $\eps$.
The coefficient of $\eps^{m-1}$ in the product of $k$ adjoint maps
is the inner sum in \eqref{eq:Yrecurrence}. Differentiation on the
left contributes $mV_m$. At $m=1$ every $j_\ell$ is zero; at $m=2$
exactly one is one, giving \eqref{eq:linearY} and
\eqref{eq:quadraticY}. Finally, \eqref{eq:Yclosed} is the
bihomogeneous projection formula \eqref{eq:allbidegrees} summed over
$p$. For a fixed total degree all sums are finite, including the
apparently infinite sum over $k$ in \eqref{eq:Yrecurrence}.
\end{proof}
Here $O(Y^2)$ means membership in the closed subspace with at least two
occurrences of $Y$, not an unspecified bound on $\norm{X}$. The linear
formula by itself is not an exact BCH formula unless a further relation
annihilates all higher-$Y$ terms.

\subsection{An independent homogeneous recursion}
\begin{theorem}\label{thm:homrec}
Let $S=X+Y$, $D=X-Y$, and
$Z(t)=\BCH(tX,tY)=\sum_{n\geq1}t^n Z_n$. Then
\begin{equation}\label{eq:compactode}
 Z'=S+\frac12[D,Z]
       +\sum_{p\geq1}b_{2p}\ad_Z^{2p}S.
\end{equation}
Consequently $Z_1=S$ and, for $n\geq2$,
\begin{equation}\label{eq:homrec}
 nZ_n=\frac12[D,Z_{n-1}]
 +\sum_{p=1}^{\lfloor(n-1)/2\rfloor}b_{2p}
 \sum_{\substack{i_1+\cdots+i_{2p}=n-1\\i_j\geq1}}
 \ad_{Z_{i_1}}\cdots\ad_{Z_{i_{2p}}}S.
\end{equation}
This recursion uniquely constructs every homogeneous term by Lie
operations, independently of the Hopf-algebra proof.
\end{theorem}
\begin{proof}
The right logarithmic derivative of $G(t)=e^{tX}e^{tY}$ is
$X+e^{t\ad_X}Y$. Also $e^{\ad_Z}Y=e^{t\ad_X}Y$, since
$e^{t\ad_Y}Y=Y$. The inverse of \eqref{eq:rightdexp} is
$z/(e^z-1)=f(z)-z$. Hence
\[
 Z'=\frac{\ad_Z}{e^{\ad_Z}-1}X+f(\ad_Z)Y.
\]
The two scalar functions have constant term $1$, opposite linear
terms $\mp z/2$, and identical even terms. Their sum applied to $X,Y$
is exactly \eqref{eq:compactode}. Coefficient comparison gives
\eqref{eq:homrec}. Its right-hand side uses only indices smaller than
$n$, proving uniqueness and the constructive assertion. Reversing
this reasoning shows that the resulting series exponentiates to
$e^{tX}e^{tY}$; thus it coincides with the logarithmic construction.
\end{proof}

\section{All displayed low-degree terms and universal identities}
\label{sec:lowdegree}
\subsection{The complete display through degree six}
In the right-nested notation \eqref{eq:rightbracket}, the terms printed
on the reference page are exactly
\begin{align}
 Z_1={}&X+Y,\label{eq:z1}\\
 Z_2={}&\tfrac12\RB{XY},\label{eq:z2}\\
 Z_3={}&\tfrac1{12}\bigl(\RB{XXY}+\RB{YYX}\bigr),\label{eq:z3}\\
 Z_4={}&-\tfrac1{24}\RB{YXXY},\label{eq:z4}\\
 Z_5={}&-\tfrac1{720}\bigl(\RB{YYYYX}+\RB{XXXXY}\bigr)
       +\tfrac1{360}\bigl(\RB{XYYYX}+\RB{YXXXY}\bigr)
       \notag\\[-2pt]
     &+\tfrac1{120}\bigl(\RB{YXYXY}+\RB{XYXYX}\bigr),
       \label{eq:z5}\\
 Z_6={}&\tfrac1{240}\RB{XYXYXY}
       +\tfrac1{720}\bigl(\RB{XYXXXY}-\RB{XXYYXY}\bigr)
       \notag\\[-2pt]
     &+\tfrac1{1440}\bigl(\RB{XYYYXY}-\RB{XXYXXY}\bigr).
       \label{eq:z6}
\end{align}
For example, $\RB{XYXXXY}=[X,[Y,[X,[X,[X,Y]]]]]$.
These displays do not select a linearly independent basis of the free
Lie algebra; Jacobi identities can produce quite different-looking
but equal answers.

\begin{proposition}
Equations \eqref{eq:z1}--\eqref{eq:z6} are the homogeneous BCH
components of degrees one through six.
\end{proposition}
\begin{proof}
We give both a derivation and finite coefficient certificates, so the
higher displayed degrees need not be accepted on the authority of a
computer program. Put $C=[X,Y]$. For $n=2$, \eqref{eq:homrec} gives
$2Z_2=\frac12[D,S]=C$. For $n=3$ it gives
$3Z_3=\frac14[D,C]$, since $\ad_S^2 S=0$. Thus
$Z_3=\frac1{12}([X,C]-[Y,C])$, equal to \eqref{eq:z3}.
For $n=4$ the same recurrence gives
\[
 4Z_4=\frac1{24}[D,[D,C]]-\frac1{24}[S,[S,C]]
       =-\frac16[Y,[X,C]].
\]
The last equality uses $[X,[Y,C]]=[Y,[X,C]]$, whose difference is
$[[X,Y],C]=[C,C]=0$. This proves \eqref{eq:z4}.

For degrees five and six, expand the right brackets in
\eqref{eq:z5}--\eqref{eq:z6} using $[A,B]=AB-BA$. The coefficient of
any word $w$ in this expansion is compared with the explicit rational
number \eqref{eq:wordcuts}. All $2^5$ and $2^6$ values of the latter,
including zeros, are printed in \cref{sec:certificate}; expanding the
respective displays gives exactly those tables. This is a finite
coefficientwise certificate of the two identities in the free
associative algebra, hence in every associative specialization.
Alternatively, substitute $Z_1,\ldots,Z_4$ in \eqref{eq:homrec} for
$n=5$, and then $Z_1,\ldots,Z_5$ for $n=6$; ordinary bracket
expansion gives the same tables. The supplied exact-arithmetic
implementation performs all three comparisons independently.
\end{proof}
An equality of noncommutative polynomials is fully determined by these
word coefficients. Thus a complete coefficient table, together with
the finite formula that computes each entry, is a proof certificate,
not a numerical approximation or a test on finitely many matrices.

\subsection{The associative display through degree four}
Expanding \eqref{eq:z1}--\eqref{eq:z4} gives the other low-order
formula on the reference page:
\begin{align}
 Z_1={}&X+Y,\notag\\
 Z_2={}&\tfrac12(XY-YX),\notag\\
 Z_3={}&\tfrac1{12}(X^2Y+XY^2-2XYX+Y^2X+YX^2-2YXY),
       \label{eq:assoclow}\\
 Z_4={}&\tfrac1{24}(X^2Y^2-2XYXY-Y^2X^2+2YXYX).\notag
\end{align}
For example,
$[Y,[X,[X,Y]]]=2XYXY-2YXYX+Y^2X^2-X^2Y^2$,
which proves the last line directly.

\subsection{Symmetry, associativity, and functoriality}
\begin{proposition}\label{prop:symmetry}
As universal formal series,
\begin{align}
 \BCH(Y,X)&=-\BCH(-X,-Y),\label{eq:swap}\\
 Z_n(Y,X)&=(-1)^{n+1}Z_n(X,Y),\label{eq:swapdegree}\\
 \BCH(\BCH(X,Y),T)&=\BCH(X,\BCH(Y,T)).\label{eq:assocbch}
\end{align}
Moreover, $\BCH(X,0)=\BCH(0,X)=X$ and
$\BCH(X,-X)=0$. Every Lie-algebra homomorphism commutes with these
series wherever formal substitution or analytic convergence is valid.
\end{proposition}
\begin{proof}
Taking inverses of $e^{-X}e^{-Y}=e^{\BCH(-X,-Y)}$ gives
$e^Ye^X=e^{-\BCH(-X,-Y)}$. Formal uniqueness of the logarithm gives
\eqref{eq:swap}, and homogeneous comparison gives
\eqref{eq:swapdegree}. Exponentiating both sides of
\eqref{eq:assocbch} gives $e^Xe^Ye^T$, so the same uniqueness proves
associativity. The identity and inverse assertions follow in the
same way. A Lie homomorphism preserves every iterated bracket, so it
preserves Dynkin's formula term by term. In a convergent analytic
setting a continuous linear homomorphism also preserves its limit;
finite-dimensional homomorphisms are automatically continuous.
\end{proof}
Let $\rho$ reverse words: $\rho(a_1\cdots a_n)=a_n\cdots a_1$.
It is an anti-automorphism and
$\rho(e^Xe^Y)=e^Ye^X$. Therefore
\begin{equation}\label{eq:reversal}
 c(\rho(w))=(-1)^{|w|+1}c(w).
\end{equation}
Both swap symmetry and reversal symmetry provide useful independent
checks on large coefficient tables.

\section{Convergence, remainder estimates, and logarithm obstructions}
\label{sec:convergence}
The formal theorem is unconditional, but its numerical evaluation is
not. This section proves the sufficient estimates stated on the
reference page, as well as several sharper distinctions needed to
interpret them. We make no claim that these elementary bounds are
optimal; convergence and continuation have a considerably richer
structure \cite{biagi}.

\subsection{Absolute convergence in a Banach algebra}
Let $\cA$ be a real or complex unital Banach algebra with a
submultiplicative norm. Write $a=\norm{X}+\norm{Y}$. The arguments
below apply equally to a finite-dimensional matrix algebra. They use
submultiplicativity, not a special choice of operator norm. In
particular the Hilbert--Schmidt matrix norm is admissible: by
Cauchy--Schwarz,
\[
 \sum_{i,j}|(AB)_{ij}|^2
 \leq\sum_{i,j}\left(\sum_k|A_{ik}|^2\right)
                    \left(\sum_k|B_{kj}|^2\right)
 =\norm{A}_{\rm HS}^2\norm{B}_{\rm HS}^2.
\]
The possible value $\norm{I}>1$ in this norm causes no difficulty,
because all the majorized logarithmic terms have positive degree.

\begin{theorem}[Two explicit convergence domains]\label{thm:convergence}
If $a<\log2$, the expanded associative series
\eqref{eq:associativeblocks} is absolutely convergent, the homogeneous
BCH series converges absolutely, and its sum $Z$ satisfies
$e^Z=e^Xe^Y$. More precisely,
\begin{equation}\label{eq:assocbound}
 \sum_{n\geq1}\norm{Z_n(X,Y)}\leq H(a),
 \qquad H(u)=-\log(2-e^u).
\end{equation}
The unreduced multiple sum of nested commutators in Dynkin's formula
is absolutely convergent whenever $a<\frac12\log2$, with bound
\begin{equation}\label{eq:dynkinbound}
 \sum_{\text{all Dynkin summands}}\norm{\text{summand}}
 \leq -\frac12\log(2-e^{2a}).
\end{equation}
\end{theorem}
\begin{proof}
First,
\[
 \norm{e^Xe^Y-I}\leq
 \sum_{r+s>0}\frac{\norm{X}^r\norm{Y}^s}{r!s!}
 =e^a-1.
\]
The sum of norms of all expanded terms of the $k$th logarithmic power
is at most $(e^a-1)^k/k$. If $a<\log2$, summing this geometric-logarithmic
majorant gives $-\log(2-e^a)$. Absolute convergence allows regrouping
by degree, and establishes \eqref{eq:assocbound}. The analytic
identity $\exp\log(I+U)=I+U$ for $\norm{U}<1$ follows by absolutely
convergent scalar-series composition, exactly as in the formal proof.
Thus the constructed $Z$ exponentiates to $e^Xe^Y$.

For the Dynkin sum, a right bracket of length $N$ has norm at most
$2^{N-1}$ times the product of its letter norms, by iterating
$\norm{[A,B]}\leq2\norm{A}\norm{B}$. In \eqref{eq:dynkin} discard
the factor $1/N\leq1$ and majorize the remaining degree-$N$ terms
using $2^{N-1}$. The $k$-block sum is then at most
$(e^{2a}-1)^k/(2k)$. Summing proves \eqref{eq:dynkinbound}.
\end{proof}
The two radii concern different ways of summing expressions. The
larger bound applies to the series of already-collected homogeneous
Lie polynomials, whereas the elementary smaller bound controls all
the unreduced terms in Dynkin's multiple sum separately.

The additional condition $\norm{Z}<\log2$ mentioned in the reference
page is not required for the logarithm \emph{constructed} by the
Mercator series. Such an additional smallness condition can be useful
when applying $\log(e^Z)=Z$ to an independently proposed logarithm,
because that reverse identity requires a branch choice. Our proof
constructs the branch and never assumes a bound on an unknown $Z$.

\begin{corollary}[Explicit truncation error]\label{cor:remainder}
If $r>1$ and $ra<\log2$, then
\begin{equation}\label{eq:remainder}
 \left\|Z-\sum_{n=1}^N Z_n\right\|
 \leq r^{-N-1}\bigl[-\log(2-e^{ra})\bigr].
\end{equation}
For a Banach Lie algebra with
$\norm{[u,v]}\leq K\norm{u}\norm{v}$, $K>0$, Dynkin's sum is
absolutely convergent if $Ka<\log2$, with total norm bounded by
$K^{-1}[-\log(2-e^{Ka})]$.
\end{corollary}
\begin{proof}
Homogeneity and \eqref{eq:assocbound} give
$\sum_{n\geq1}r^n\norm{Z_n}\leq H(ra)$. Every term with $n>N$
therefore contributes at most $r^{-N-1}$ times its weighted norm.
For the Lie-algebra assertion replace the factor $2^{N-1}$ in the
proof of \eqref{eq:dynkinbound} by $K^{N-1}$. If the bracket is zero,
BCH simply terminates at $X+Y$ and no positive $K$ is needed.
\end{proof}
Every finite-dimensional Lie algebra admits such a $K$, because a
bilinear map is bounded on the product of unit spheres. This supplies
an elementary local convergence proof without requiring a matrix
representation of the Lie algebra.

\subsection{The trace identity and its branch restriction}
\begin{proposition}\label{prop:trace}
For finite matrices on the local BCH branch,
\begin{equation}\label{eq:trace}
 \tr\log(e^Xe^Y)=\tr X+\tr Y.
\end{equation}
For an arbitrary matrix logarithm $L$ of $e^Xe^Y$ over $\C$, the
unconditional assertion is only
\begin{equation}\label{eq:tracebranch}
 \tr L-\tr X-\tr Y\in 2\pi i\Z.
\end{equation}
\end{proposition}
\begin{proof}
The equality $\tr(AB)=\tr(BA)$ follows by interchanging the two finite
summation indices. Therefore every commutator has zero trace, and
all $Z_n$ for $n\geq2$ have zero trace. Continuity of trace proves
\eqref{eq:trace} wherever the BCH series converges.
To prove the second statement, triangularize a complex matrix $A$:
an eigenvector exists, and induction on the quotient dimension gives
an upper triangular matrix with diagonal $\lambda_1,\ldots,\lambda_d$.
Its exponential is upper triangular with diagonal
$e^{\lambda_1},\ldots,e^{\lambda_d}$. Hence
$\det e^A=e^{\tr A}$. Apply this identity to $L,X,Y$ and use
multiplicativity of determinant. The scalar identity $e^z=1$ holds
exactly when $z\in2\pi i\Z$, giving \eqref{eq:tracebranch}.
\end{proof}
For example, take scalar $X=Y=3\pi i/4$. Their sum is $3\pi i/2$,
but the principal logarithm of $e^Xe^Y=-i$ is $-\pi i/2$. Thus the
word ``logarithm'' cannot be left branch-independent in
\eqref{eq:trace}.

\subsection{A complete proof of the \texorpdfstring{$2\times2$}{2 by 2} obstruction}
\begin{theorem}\label{thm:obstruction}
For
\begin{equation}\label{eq:obstructionXY}
 X=\begin{pmatrix}0&i\pi\\i\pi&0\end{pmatrix},\qquad
 Y=\begin{pmatrix}0&1\\0&0\end{pmatrix},
\end{equation}
the product is
\begin{equation}\label{eq:obstructionproduct}
 e^Xe^Y=\begin{pmatrix}-1&-1\\0&-1\end{pmatrix}.
\end{equation}
It has no traceless complex logarithm. Its homogeneous BCH series
therefore diverges, although the product does have logarithms in
$\mathfrak{gl}_2(\C)$.
\end{theorem}
\begin{proof}
Since $X^2=-\pi^2I$, the even and odd exponential series give
$e^X=-I$. Since $Y^2=0$, $e^Y=I+Y$, proving
\eqref{eq:obstructionproduct}.
Any logarithm $L$ commutes with its own exponential, hence with $Y$.
Writing out $LY=YL$ shows that the centralizer of $Y$ consists of
matrices $aI+bY$. For such a matrix,
$e^{aI+bY}=e^a(I+bY)$. Equality with $-I-Y$ forces
$e^a=-1$ and $b=1$. Thus \emph{all} logarithms are
\begin{equation}\label{eq:allobstructionlogs}
 L=(2k+1)\pi iI+Y,\qquad k\in\Z,
\end{equation}
and each has nonzero trace $2(2k+1)\pi i$.

It remains to justify that even conditional convergence of BCH would
contradict this fact. Suppose $\sum_{n\geq1}Z_n(X,Y)=S$ converged.
Its matrix power series $Z(t)=\sum_{n\geq1}t^nZ_n$ would converge
for $|t|<1$, and $e^{Z(t)}=e^{tX}e^{tY}$ there: the identity holds
near zero by the local theorem and then throughout the disk by
analytic uniqueness, entry by entry. Abel's limit theorem gives
$Z(t)\to S$ as real $t\uparrow1$. To see that theorem directly,
let $S_m=\sum_{n=1}^m Z_n$. Summation by parts gives
$Z(t)=(1-t)\sum_{m\geq1}t^m S_m$; split the sum into a fixed finite
head and a tail where $S_m$ is uniformly close to $S$. The head tends
to zero and the tail tends to $S$. Continuity of the exponential
now yields $e^S=e^Xe^Y$. But every $Z_n$ is traceless, so $S$ is
traceless, contradicting \eqref{eq:allobstructionlogs}.
\end{proof}
This example simultaneously distinguishes existence of a logarithm,
existence of a logarithm in the specified Lie algebra, and convergence
of the BCH series. They are genuinely different questions.

\section{Lie groups, local multiplication, and integration}
\label{sec:localgroups}
\subsection{Matrix groups and the intrinsic differential}
For a matrix Lie group $G$, its Lie algebra is the tangent space at
the identity. The left-invariant field associated with a tangent
matrix $X$ is $X^L(g)=gX$. Its bracket with $Y^L$ is
$g(XY-YX)$, as follows by differentiating the two vector fields.
Thus the induced Lie bracket is the matrix commutator. The integral
curve through $I$ of $X^L$ solves $g'=gX$, $g(0)=I$, and its
solution is the ordinary matrix exponential $e^{tX}$.

The relevant assertions do not require a matrix embedding. For an
arbitrary finite-dimensional Lie group define its left
Maurer--Cartan form by $\theta_g=(L_{g^{-1}})_*$. On left-invariant
fields it takes constant values. The formula for the exterior
derivative of a one-form consequently gives
\begin{equation}\label{eq:MC}
 d\theta+\tfrac12[\theta,\theta]=0.
\end{equation}
Indeed, evaluated on $U^L,V^L$, the first term is $-[U,V]$ and the
second is $[U,V]$; left-invariant fields span each tangent space.

For an intrinsic proof of \eqref{eq:leftdexp}, consider
$g(t,u)=\exp(t(X+uH))$ and set
$\alpha=\theta(\partial_tg)$, $\beta=\theta(\partial_ug)$.
Here $\alpha=X+uH$. Pulling back \eqref{eq:MC} gives
$\partial_t\beta-\partial_u\alpha=-[\alpha,\beta]$.
At $u=0$, therefore,
\[
 \partial_t\beta=H-\ad_X\beta,\qquad\beta(0)=0.
\]
Solving this linear equation gives
$\beta(1)=\int_0^1e^{-s\ad_X}H\,ds=\phi_-(\ad_X)H$.
This is exactly the left logarithmic differential of the exponential.
In particular $(d\exp)_0=\id$.

\begin{theorem}[Local BCH multiplication]\label{thm:localgroup}
The exponential map of a finite-dimensional real or complex Lie
group is a local diffeomorphism at zero. In its local inverse chart,
\begin{equation}\label{eq:localgroup}
 \exp X\,\exp Y=\exp\BCH(X,Y)
\end{equation}
for sufficiently small $X,Y$. Hence the local multiplication law is
determined entirely by the Lie bracket.
\end{theorem}
\begin{proof}
The inverse function theorem applies because $(d\exp)_0=\id$.
Choose a norm and a bracket bound $K$ as in
\cref{cor:remainder}; the universal BCH series then converges in a
small ball in the Lie algebra. Its termwise differential equation
\eqref{eq:bchode} is valid in a smaller ball, by uniform convergence
on compact subballs, or directly by the convergent majorant with a
slightly larger scaling parameter. The intrinsic differential just
proved shows that $\exp Z(t)$ has left logarithmic derivative $Y$
and starts at $\exp X$. Uniqueness for the corresponding smooth
ordinary differential equation gives
$\exp Z(t)=\exp X\exp(tY)$. At $t=1$ this is
\eqref{eq:localgroup}. Local injectivity of the exponential selects
this logarithm uniquely.
\end{proof}
The local assertion is not a global classification: the additive
Lie group $\R$ and the circle $\{e^{it}:t\in\R\}$ have the same
one-dimensional abelian Lie algebra, but one is noncompact and the
other compact, so they cannot be isomorphic as Lie groups.

\subsection{Homomorphisms and representations}
\begin{theorem}[Local Lie correspondence]\label{thm:localcorrespondence}
A Lie-algebra homomorphism $A:\g\to\h$ defines a unique germ at the
identity of a Lie-group homomorphism by
\begin{equation}\label{eq:localhom}
 F(\exp_G X)=\exp_H(AX).
\end{equation}
Conversely, the differential at the identity of a smooth local
group homomorphism preserves Lie brackets.
\end{theorem}
\begin{proof}
On exponential neighborhoods formula \eqref{eq:localhom} defines a
smooth map. Naturality of BCH and \cref{thm:localgroup} show that it
preserves products whenever the inputs and their product remain in
the chosen neighborhoods. For the converse, differentiating
$F(gh)=F(g)F(h)$ shows that $F$ relates the left-invariant field of
$X$ to that of $(dF)_eX$. The chain rule for vector-field brackets
then gives $(dF)_e[X,Y]=[(dF)_eX,(dF)_eY]$. Finally, a local
homomorphism sends the integral curve $\exp_G(tX)$ to
$\exp_H(t(dF)_eX)$, proving uniqueness of its germ.
\end{proof}

\begin{theorem}[Integration from a simply connected group]
\label{thm:integration}
If $G$ is connected and simply connected and $H$ is a Lie group,
every Lie-algebra homomorphism $A:\g\to\h$ integrates to a unique
smooth homomorphism $F:G\to H$.
\end{theorem}
\begin{proof}
Start with the local homomorphism $F_0$ of
\cref{thm:localcorrespondence}, and shrink its identity neighborhood
so that sufficiently short products are covered by the local
multiplication rule. Given a path $\gamma$ from $e$ to $g$, choose a
partition fine enough that all increments
$\gamma(t_j)^{-1}\gamma(t_{j+1})$ lie in this neighborhood, and set
\[
 F_\gamma(g)=\prod_j
 F_0\bigl(\gamma(t_j)^{-1}\gamma(t_{j+1})\bigr).
\]
Inserting a partition point does not change the product, because
the old increment is the product of the two new increments and
$F_0$ preserves that small product. A common refinement therefore
shows independence of a sufficiently fine partition.

For homotopy invariance, cover the compact parameter square of an
endpoint-fixed homotopy by small rectangles on which all group
values differ by elements of a still smaller identity neighborhood.
Uniform continuity provides a finite rectangular subdivision with
this property. Around each small rectangle, the four increments
multiply to the identity. By further shrinking at the outset,
all partial products around such a rectangle are in the domain
where $F_0$ is multiplicative. Their images therefore multiply to
the identity as well. Replacing one edge path of a rectangle by the
other preserves its product. Sweeping across the grid proves equality
for its two boundary paths; the endpoint edges contribute identities.
The partition argument handles the actual boundary curves rather
than just their grid descriptions. Thus $F_\gamma$ depends only on
the endpoint-fixed homotopy class of $\gamma$.

Simple connectedness makes it depend only on $g$; call it $F(g)$.
Concatenate a path to $g$ with the left translate by $g$ of a path
to $h$. The increments of the second path are the original
increments of the path to $h$, so $F(gh)=F(g)F(h)$.
For $u$ close to $e$, appending a short path from $g$ to $gu$ gives
$F(gu)=F(g)F_0(u)$, proving smoothness. Any other homomorphism with
differential $A$ agrees with $F_0$ near the identity and hence on
finite products of such a neighborhood. These products form an
open subgroup whose cosets are open; connectedness makes it all
of $G$. This proves uniqueness.
\end{proof}
Taking $H=\GL(W)$ proves the corresponding integration statement for
finite-dimensional representations. Without simple connectedness,
the path construction can have a nontrivial value around a loop;
this is the precise additional global obstruction absent from the
local BCH law.

\subsection{Nilpotent algebras and a global polynomial group law}
\begin{proposition}\label{prop:nilpotent}
If a Lie algebra is nilpotent of class at most $c$, every BCH term
of degree greater than $c$ vanishes. On its underlying vector space,
\begin{equation}\label{eq:nilpotentlaw}
 x*y=\sum_{n=1}^{c}Z_n(x,y)
\end{equation}
defines a Lie group with identity $0$, inverse $-x$, original Lie
bracket, and exponential map equal to the identity of the vector
space. This group is connected and simply connected.
\end{proposition}
\begin{proof}
For the lower central series $\g_1=\g$,
$\g_{j+1}=[\g,\g_j]$, induction using Jacobi shows
$[\g_p,\g_q]\subseteq\g_{p+q}$. Every Lie monomial of length $n$
therefore lies in $\g_n$ and vanishes for $n>c$.
The formal associativity, identity, and inverse identities of
\cref{prop:symmetry} now become finite polynomial identities on all
of $\g$. Thus \eqref{eq:nilpotentlaw} is a smooth global group law.
Its quadratic term is $[x,y]/2$; differentiating the group
commutator $(tx)*(sy)*(-tx)*(-sy)$ at $t=s=0$ gives $[x,y]$.
Alternatively this follows by direct differentiation of its
left-invariant fields. Since all brackets of multiples of a fixed
$x$ vanish, $(tx)*(sx)=(t+s)x$. These are exactly the one-parameter
groups, so the exponential is $x\mapsto x$. The underlying manifold
is a vector space and is therefore connected and simply connected.
\end{proof}
If a given connected simply connected Lie group has this nilpotent
Lie algebra, integrating the identity Lie-algebra maps in both
directions by \cref{thm:integration} gives inverse group
isomorphisms. Its exponential is consequently a global
diffeomorphism. For any other connected group with the same Lie
algebra, integration from the polynomial group still shows that its
exponentials satisfy the same finite BCH multiplication formula.

\section{Exact special cases and their analytic limits}
\label{sec:special}
\subsection{Commuting inputs and a central commutator}
\begin{theorem}[Central-commutator BCH]\label{thm:central}
Suppose $C=[X,Y]$ satisfies $[X,C]=[Y,C]=0$. Then
\begin{align}
 e^Xe^Y&=e^{X+Y+C/2},\label{eq:centralbch}\\
 e^{X+Y}&=e^Xe^Ye^{-C/2}
        =e^{X/2}e^Ye^{X/2},\label{eq:centralsym}\\
 e^Xe^Ye^{-X}e^{-Y}&=e^C.\label{eq:groupcomm}
\end{align}
These are formal identities, global identities for bounded elements
of a Banach algebra, and global identities for elements of a
finite-dimensional Lie group satisfying these bracket relations.
In particular, $[X,Y]=0$ gives $e^Xe^Y=e^{X+Y}$.
\end{theorem}
\begin{proof}
Hadamard's identity truncates to
$e^{sX}Ye^{-sX}=Y+sC$. Thus the right logarithmic derivative of
$G(s)=e^{sX}e^{sY}$ is $X+Y+sC$. The function
$E(s)=e^{s(X+Y)+s^2C/2}$ has the same derivative, since $X+Y$ and
$C$ commute. Both start at the identity, proving
\eqref{eq:centralbch} at $s=1$ by uniqueness. The intrinsic
logarithmic derivative gives the identical proof in a Lie group.
Moving the central exponential to the other side gives the first
part of \eqref{eq:centralsym}. Applying \eqref{eq:centralbch}
twice to $X/2,Y,X/2$ cancels the two half-commutator corrections
and gives its symmetric part. Finally, conjugation sends $Y$ to
$Y+C$, so $e^Xe^Ye^{-X}e^{-Y}=e^{Y+C}e^{-Y}=e^C$.
\end{proof}
This proof is global and does not infer convergence from the
formal truncation alone. For unbounded operators a common invariant
domain and integrability assumptions are still necessary; see
\cref{sec:quantum}.

If all pairwise commutators of $A_1,\ldots,A_m$ commute with every
$A_j$, induction on $m$ gives the full finite-factor identity
\begin{equation}\label{eq:centralmulti}
 \prod_{j=1}^m e^{A_j}
 =\exp\left(\sum_{j=1}^m A_j+
             \frac12\sum_{j<k}[A_j,A_k]\right).
\end{equation}
Indeed, the commutator of the exponent for the first $m-1$ factors
with $A_m$ is just $\sum_{j<m}[A_j,A_m]$, since its central part
commutes with $A_m$. The two-factor formula completes the induction.
For a basis $P,Q,C$ with $[P,Q]=C$ and all other brackets zero,
\eqref{eq:nilpotentlaw} is the Heisenberg law
\begin{equation}\label{eq:heisenberglaw}
 (x,y,z)*(x',y',z')=
 (x+x',y+y',z+z'+\tfrac12(xy'-yx')).
\end{equation}
This also proves directly that the Heisenberg algebra is two-step
nilpotent and that all higher BCH terms vanish.

\subsection{An adjoint eigenvector and exact resummation}
\begin{theorem}\label{thm:eigen}
Suppose $[X,Y]=sY$, where $s$ is a scalar. Then
\begin{align}
 e^Xe^Y&=e^{e^sY}e^X,\label{eq:eigenbraid}\\
 e^{X+cY}&=e^X e^{c\phi_-(s)Y}\qquad(c\text{ scalar}).
 \label{eq:eigenfactor}
\end{align}
If $s\notin2\pi i\Z\setminus\{0\}$, this implies
\begin{equation}\label{eq:eigenbch}
 e^Xe^Y=\exp\left(X+\frac{s}{1-e^{-s}}Y\right),
\end{equation}
with coefficient $1$ at $s=0$. These exponential identities are
global in the same bounded-algebra or finite-dimensional-group
settings as \cref{thm:central}. The local BCH Taylor series in this
case is $X+f(s)Y$ with $f$ expanded at zero; resummation and Taylor
convergence are distinct assertions.
\end{theorem}
\begin{proof}
The relation implies $\ad_X^nY=s^nY$, so Hadamard's formula gives
$e^XYe^{-X}=e^sY$, proving \eqref{eq:eigenbraid} by conjugating the
exponential. For \eqref{eq:eigenfactor}, put
\[
 q(t)=c\int_0^t e^{-su}\,du,
 \qquad G(t)=e^{tX}e^{q(t)Y}.
\]
Its right logarithmic derivative is
$X+q'(t)e^{ts}Y=X+cY$. Since $G(0)=I$, it equals
$e^{t(X+cY)}$. At $t=1$, $q(1)=c\phi_-(s)$, proving the
factorization. Choosing $c=1/\phi_-(s)=f(s)$ is possible exactly
under the stated nonresonance condition and proves
\eqref{eq:eigenbch}.

For the formal BCH assertion, the Lie algebra generated by $X,Y$
is contained in their span, and its ideal spanned by $Y$ is abelian.
Every bracket monomial containing at least two copies of $Y$
therefore vanishes: induct on a binary bracket expression, using
that a bracket containing one $Y$ is a multiple of $Y$, while an
all-$X$ bracket of length at least two is zero. Consequently only
\eqref{eq:linearY} remains. Applying $\ad_X^nY=s^nY$ gives the
Taylor expansion of $f(s)$.
\end{proof}
At a nonzero resonance the restriction is genuine. Take
\[
 X=\begin{pmatrix}2\pi i&0\\0&0\end{pmatrix},\qquad
 Y=\begin{pmatrix}0&1\\0&0\end{pmatrix}.
\]
Then $[X,Y]=2\pi iY$, $e^Xe^Y=I+Y$, whereas
\eqref{eq:eigenfactor} gives $e^{X+cY}=I$ for every $c$.
There is nevertheless a logarithm of the product, namely $Y$.
Thus the failure concerns the particular form $X+cY$, not the
existence of any logarithm.

There is also a nonresonant distinction. If $Y\ne0$ and $s\ne0$,
\begin{equation}\label{eq:eigenscaled}
 \BCH(tX,tY)=tX+t f(ts)Y
\end{equation}
as a Taylor series at $t=0$, and its radius is exactly
$2\pi/|s|$, because the poles at $ts=\pm2\pi i$ have nonzero
residues in the $Y$ direction. Thus real $s>2\pi$ can make this
series divergent at $t=1$, although the resummed exponential
identity \eqref{eq:eigenbch} remains valid. The conjugation law
\eqref{eq:eigenbraid}, being an entire exponential identity, has
no such resonance exclusion at all.

\section{Zassenhaus factorization to every order}
\label{sec:zassenhaus}
The complementary problem is to factor the exponential of a sum
into exponentials of homogeneous Lie corrections. We first solve
this in the degree completion, then prove a concrete analytic
convergence theorem. Recursive constructions in this direction are
also central in the computational literature \cite{casas}.

\subsection{Existence, uniqueness, and primitive corrections}
\begin{theorem}[Formal Zassenhaus product]\label{thm:zassenhaus}
There is a unique sequence of homogeneous Lie polynomials $C_n$ of
degree $n$, $n\geq2$, for which
\begin{equation}\label{eq:zassenhaus}
 e^{t(X+Y)}=e^{tX}e^{tY}
             e^{t^2C_2}e^{t^3C_3}e^{t^4C_4}\cdots.
\end{equation}
The infinite product is read from left to right in increasing degree
and is defined coefficientwise in $t$.
\end{theorem}
\begin{proof}
Set
\begin{equation}\label{eq:residualstart}
 R_1(t)=e^{-tY}e^{-tX}e^{t(X+Y)}=I+O(t^2).
\end{equation}
Inductively, if $R_{n-1}=I+O(t^n)$, define
\begin{equation}\label{eq:residualalgorithm}
 C_n=[t^n]R_{n-1}(t),\qquad
 R_n(t)=e^{-t^nC_n}R_{n-1}(t).
\end{equation}
Then $R_n=I+O(t^{n+1})$. Each residual is group-like: this is true
initially by \cref{lem:grouplike}, and remains true after removing
the exponential of a primitive coefficient. To verify the induction,
compare degree $n$ in
$\Delta R_{n-1}=R_{n-1}\otimes R_{n-1}$. Its first nonconstant
coefficient satisfies
$\Delta C_n=C_n\otimes1+1\otimes C_n$. By
\cref{thm:friedrichs} it is a Lie polynomial. Its degree is $n$,
since every input has matching degree in $t$ and the letters.

Undoing the residuals gives
$R_1=e^{t^2C_2}\cdots e^{t^nC_n}R_n$.
As $n\to\infty$, $R_n\to I$ coefficientwise, proving
\eqref{eq:zassenhaus}. In any proposed factorization, the first
unmatched coefficient must be $[t^n]R_{n-1}$, proving uniqueness.
\end{proof}
In the alternative display
$e^{-tX}e^{t(X+Y)}=\prod_{n\geq1}e^{t^nC_n}$, the convention
\emph{must} be $C_1=Y$. The notation would otherwise omit a
necessary first-degree term.

\subsection{The displayed exponents and a differential algorithm}
The first corrections are
\begin{align}
 C_2&=-\tfrac12[X,Y],\label{eq:c2}\\
 C_3&=\tfrac16[X,[X,Y]]+\tfrac13[Y,[X,Y]],\label{eq:c3}\\
 C_4&=-\tfrac1{24}\bigl([[[X,Y],X],X]
           +3[[[X,Y],X],Y]+3[[[X,Y],Y],Y]\bigr)\notag\\
     &=-\tfrac1{24}\RB{XXXY}
       -\tfrac18\RB{YXXY}-\tfrac18\RB{YYXY}.
       \label{eq:c4}
\end{align}
Both versions of $C_4$ agree by antisymmetry; for example
$[[[X,Y],X],Y]=[Y,[X,[X,Y]]]$.

For a complete algorithm using only brackets, set
$F_n=R_n'R_n^{-1}=\sum_{k\geq0}t^k f_{n,k}$.
Direct differentiation of \eqref{eq:residualstart}, using Hadamard's
formula, yields
\begin{equation}\label{eq:F1}
 F_1=e^{-t\ad_Y}(e^{-t\ad_X}-I)Y
 =\sum_{k\geq1}(-t)^k
   \sum_{j=1}^{k}\frac{\ad_Y^{k-j}\ad_X^jY}{(k-j)!j!}.
\end{equation}
Indeed the right logarithmic derivatives of the three factors give
$-Y-e^{-t\ad_Y}X+e^{-t\ad_Y}e^{-t\ad_X}(X+Y)$;
the $X$ terms cancel, and $e^{-t\ad_Y}Y=Y$.
Differentiating \eqref{eq:residualalgorithm} gives, for $n\geq2$,
\begin{align}
 F_n&=e^{-t^n\ad_{C_n}}F_{n-1}-nt^{n-1}C_n,
       \label{eq:Fn}\\
 C_n&=\frac1n f_{n-1,n-1},\label{eq:CfromF}\\
 f_{n,k}&=\sum_{j=0}^{\lfloor k/n\rfloor}
       \frac{(-1)^j}{j!}\ad_{C_n}^j f_{n-1,k-nj}
       -n C_n\,\delta_{k,n-1}.\label{eq:Fcoeff}
\end{align}
Equation \eqref{eq:CfromF} follows from
$R_{n-1}=I+t^nC_n+O(t^{n+1})$. These formulas are valid at every
order and constitute a second construction of the same unique
factorization.

For an explicit derivation of \eqref{eq:c2}--\eqref{eq:c4}, write
$C=[X,Y]$. Formula \eqref{eq:F1} begins
\[
 F_1=-tC+t^2\bigl(\tfrac12[X,C]+[Y,C]\bigr)
 -t^3\bigl(\tfrac16[X,[X,C]]
            +\tfrac12[Y,[X,C]]+\tfrac12[Y,[Y,C]]\bigr)
 +O(t^4).
\]
Dividing the $t$ coefficient by $2$ gives $C_2$. Passing to $F_2$
removes that coefficient; its possible degree-three modification is
$-[C_2,-C]=0$. Hence its degree-two coefficient divided by $3$
gives $C_3$, and its degree-three coefficient divided by $4$ gives
$C_4$. Removing $C_3$ does not affect the latter coefficient.
This proves every displayed exponent without leaving implicit a
higher-order approximation that could change it.

\subsection{A finite nonrecursive formula using weighted trees}
The previous algorithms determine all coefficients recursively.
For a genuinely finite nonrecursive expression at every fixed
order, first define the explicit associative polynomials
\begin{equation}\label{eq:Hn}
 H_n=[t^n]R_1(t)=
 \sum_{a+b+c=n}\frac{(-1)^{a+b}}{a!b!c!}
       Y^aX^b(X+Y)^c.
\end{equation}
Here $H_0=I$ and $H_1=0$; each $(X+Y)^c$ is the sum of all $2^c$
words of length $c$. Coefficient extraction from the ordered product
for $R_1$ gives
\begin{equation}\label{eq:forwardC}
 C_n=H_n-
 \sum_{\substack{k_2,\ldots,k_{n-1}\geq0\\
                 \sum_{j=2}^{n-1}jk_j=n}}
 \frac{C_2^{k_2}\cdots C_{n-1}^{k_{n-1}}}
      {k_2!\cdots k_{n-1}!}.
\end{equation}
The sum is empty for $n=2,3$. Every product in it has at least two
factors. Importantly, its order is increasing in the degree of its
factors: these polynomials need not commute.

\begin{definition}
Let $\cT_n$ consist of finite plane rooted trees with an integer
weight at each vertex. The root has weight $n$, every weight is at
least $2$, and each internal vertex has at least two children.
The child weights are weakly increasing from left to right and sum
to the parent's weight. Leaves have no further restriction.
For an internal vertex $v$, let $m_v(j)$ be the number of its
children of weight $j$. Define
\begin{equation}\label{eq:treeweight}
 \omega(T)=(-1)^{\#\{\text{internal vertices}\}}
           \prod_{v\ \mathrm{internal}}\prod_{j\geq2}\frac1{m_v(j)!},
 \qquad
 H(T)=\prod_{\ell\ \mathrm{leaf,\ left\ to\ right}}H_{\mathrm{wt}(\ell)}.
\end{equation}
Subtrees of equal root weight are still ordered; interchanging
unequal such subtrees can change $H(T)$.
\end{definition}

\begin{theorem}[All-terms Zassenhaus tree formula]\label{thm:trees}
For every $n\geq2$,
\begin{equation}\label{eq:treeformula}
 \boxed{\displaystyle C_n=\sum_{T\in\cT_n}\omega(T)H(T).}
\end{equation}
This is a finite formula involving only factorials, sums, and
products of $X,Y$. An explicitly nested-commutator version is
\begin{equation}\label{eq:treelie}
 C_n=\frac1n\sum_{T\in\cT_n}\omega(T)
       \sum_{w\in\{X,Y\}^n}([w]H(T))\RB{w}.
\end{equation}
\end{theorem}
\begin{proof}
All leaf weights sum to $n$, so there are at most $\lfloor n/2\rfloor$
leaves. With at least two children at every internal vertex there
are fewer internal vertices than leaves. Thus only finitely many
tree shapes and weight assignments occur.
The tree consisting only of its root contributes $H_n$. For any
other tree, let the root's child weights be the nondecreasing list
$i_1,\ldots,i_k$. Removing the root leaves an independently chosen
tree in each $\cT_{i_j}$. Its root contributes the negative sign
and the factor $\prod_jm(j)!^{-1}$. Therefore the right-hand side
of \eqref{eq:treeformula}, grouped by root children, satisfies
exactly \eqref{eq:forwardC}, with initial values $H_2,H_3$.
Induction on $n$ proves the formula. Since $C_n$ is primitive,
applying the right Dynkin projection $D/n$ to it leaves it fixed.
Applying that projection to its word expansion gives
\eqref{eq:treelie}.
\end{proof}
For orientation, the first nontrivial instances are
\begin{align}
 C_2&=H_2,\quad C_3=H_3,\quad C_4=H_4-\tfrac12H_2^2,\notag\\
 C_5&=H_5-H_2H_3,\notag\\
 C_6&=H_6-H_2H_4-\tfrac12H_3^2+\tfrac13H_2^3.
 \label{eq:CinH}
\end{align}
The last coefficient is $1/2-1/6=1/3$: one contribution comes
from substituting $C_4=H_4-H_2^2/2$ in $-C_2C_4$, and the other
from $-C_2^3/6$. This example illustrates why an unordered product
would lose information at higher degrees.

\subsection{A quantitative convergence theorem for the product}
\begin{theorem}\label{thm:Zconv}
For bounded $X,Y$ in a Banach algebra put
$a=\norm{X}+\norm{Y}$. If $r>0$ satisfies
\begin{equation}\label{eq:Zcriterion}
 h(r):=e^{2ar}-1-2ar<2\log2-1,
\end{equation}
then $\sum_{n\geq2}r^n\norm{C_n}<\infty$ and the Zassenhaus product
converges in norm for $|t|\leq r$, to $e^{t(X+Y)}$.
A particularly simple sufficient condition is $a|t|\leq1/4$.
\end{theorem}
\begin{proof}
The three exponential factors defining $H_n$ give
$\norm{H_n}\leq(2a)^n/n!$ for $n\geq2$.
Introduce a scalar formal series $c(z)$ with nonnegative
coefficients by
\begin{equation}\label{eq:majorantimplicit}
 c=h+(e^c-1-c),\qquad h(z)=e^{2az}-1-2az.
\end{equation}
The coefficient of degree $n$ in $e^c-1-c$ depends only on
coefficients of $c$ of degrees less than $n$, since $c$ starts at
degree $2$. Thus this defines a unique nonnegative formal series.
The triangle inequality in \eqref{eq:forwardC}, and the fact that
scalar factors commute, prove inductively that
$\norm{C_n}\leq[z^n]c(z)$: the scalar product coefficients are
exactly those of $e^{c(z)}$ after removing the constant and linear
terms in $c$.

To prove convergence rather than just formal majorization, choose
$q\in(0,\log2)$ with $h(r)=2q+1-e^q$. Such a $q$ exists under
\eqref{eq:Zcriterion}, since $2q+1-e^q$ is strictly increasing
from $0$ to $2\log2-1$ on that interval. The case $a=0$ is
trivial. Iterate $c^{(0)}=0$ and
$c^{(m+1)}=h+e^{c^{(m)}}-1-c^{(m)}$. Its coefficients increase
monotonically and stabilize at each fixed degree to those of $c$.
At $z=r$ the scalar values stay at most $q$, by induction and the
monotonicity of $e^u-1-u$ for $u\geq0$. Monotone convergence of
nonnegative series gives $\sum_n[z^n]c\,r^n\leq q$.

It follows that $\sum\norm{t^nC_n}\leq q$ for $|t|\leq r$.
Successive products of exponentials are Cauchy in norm: a tail
product differs from $I$ by at most
$\exp(\sum_{n>N}\norm{t^nC_n})-1$, and the norm of the preceding
product is uniformly bounded by the exponential of the total sum
(up to the harmless norm of $I$). Absolute convergence permits
coefficientwise multiplication, so the formal identity
\eqref{eq:zassenhaus} proves equality of the analytic products for
$|t|<r$, and continuity proves equality on the boundary.
Finally $a|t|\leq1/4$ implies
$h(|t|)\leq e^{1/2}-3/2<2\log2-1$, as is seen, for example,
from the convergent scalar series or elementary numerical bounds
$e^{1/2}<1.649$ and $\log2>0.693$.
\end{proof}
This is an elementary sufficient criterion, not a claim about the
maximal domain or the best constant in the literature.

If $[X,Y]=sY$, the factorization \eqref{eq:eigenfactor} at scale $t$
gives a particularly simple all-orders answer:
\begin{equation}\label{eq:eigenZassenhaus}
 C_n=\frac{(-1)^{n-1}s^{n-1}}{n!}\,Y\qquad(n\geq2).
\end{equation}
Indeed, $t\phi_-(ts)Y=tY+
\sum_{n\geq2}(-1)^{n-1}t^ns^{n-1}Y/n!$, and all these exponents
commute. This Zassenhaus product converges for every $t$, even in
examples where the BCH Taylor series has finite radius. If the
commutator is central instead, only $C_2=-[X,Y]/2$ survives.

\section{Trotter and Suzuki product formulas with full error series}
\label{sec:splitting}
The reference page mentions the Suzuki--Trotter decomposition as a
consequence of Zassenhaus. Here we give the bounded-operator results
and their all-order logarithmic expansions. They do not assert an
unrestricted product theorem for unbounded semigroup generators.
The classical sources treat broader settings \cite{trotter,suzuki}.

\subsection{The ordinary product formula}
\begin{theorem}\label{thm:trotter}
For bounded $X,Y$ and every fixed real or complex $t$,
\begin{equation}\label{eq:trotter}
 \lim_{n\to\infty}(e^{tX/n}e^{tY/n})^n=e^{t(X+Y)}
\end{equation}
in norm, locally uniformly in $t$. For sufficiently large $n$ its
complete logarithmic expansion is
\begin{equation}\label{eq:trotterall}
 (e^{tX/n}e^{tY/n})^n=
 \exp\left(t(X+Y)+\sum_{k\geq2}t^k n^{1-k}Z_k(X,Y)\right).
\end{equation}
The convergence error is $O(n^{-1})$, uniformly on each compact
$t$-set.
\end{theorem}
\begin{proof}
When $|t|a/n<\log2$, the one-step logarithm is the convergent BCH
series. Since powers of a single exponential satisfy
$(e^A)^n=e^{nA}$, multiplying its logarithm by $n$ gives
\eqref{eq:trotterall}. The sum of terms starting at $k=2$ is
$O(n^{-1})$ uniformly for bounded $t$, by the convergence majorant
on any fixed slightly larger disk in the variable $t/n$.
For bounded $A,B$, differentiation of
$e^{(1-s)A}e^{sB}$ gives the identity
\[
 e^B-e^A=\int_0^1e^{(1-s)A}(B-A)e^{sB}\,ds.
\]
Consequently $\norm{e^B-e^A}\leq
\norm{B-A}e^{\max(\norm{A},\norm{B})}$ for a normalized norm;
the same estimate with a fixed harmless multiplicative constant
holds in the other admissible norms. Applying it to the exponents
proves the stated limit and error.
\end{proof}
Every coefficient of the error series is explicitly supplied by
\eqref{eq:wordcuts} or Dynkin's formula; it is not merely a formal
symbol for an unknown remainder.

\subsection{The symmetric product and arbitrary even order}
Define
\begin{equation}\label{eq:S2}
 S_2(t)=e^{tX/2}e^{tY}e^{tX/2}.
\end{equation}
Since $S_2(-t)=S_2(t)^{-1}$, its local logarithm is odd:
\begin{equation}\label{eq:symmetriclog}
 \log S_2(t)=t(X+Y)+\sum_{j\geq1}t^{2j+1}E_{2j+1}.
\end{equation}
A direct substitution into BCH through degree three gives
\begin{equation}\label{eq:symmetriccubic}
 E_3=-\frac1{24}[X,[X,Y]]-\frac1{12}[Y,[X,Y]].
\end{equation}
For clarity, one can obtain this by first writing
$A=\BCH(tX/2,tY)=tX/2+tY+t^2[X,Y]/4+
 t^3([X,[X,Y]]/48-[Y,[X,Y]]/24)+O(t^4)$,
and then expanding $\BCH(A,tX/2)$ using
\eqref{eq:z2}--\eqref{eq:z3}. The quadratic terms cancel and the
cubic terms combine to \eqref{eq:symmetriccubic}; all even terms
vanish by the inverse symmetry. The multifactor version of
\cref{thm:wordcuts} gives every $E_{2j+1}$ as a finite cut sum.
It follows as above that $S_2(t/n)^n$ approximates $e^{t(X+Y)}$
with norm error $O(n^{-2})$ and has full logarithmic expansion
$t(X+Y)+\sum_{j\geq1}t^{2j+1}n^{-2j}E_{2j+1}$.

\begin{theorem}[Suzuki's even-order recursion]\label{thm:suzuki}
For $k\geq2$, set
\begin{equation}\label{eq:suzukip}
 p_k=\frac1{4-4^{1/(2k-1)}},
\end{equation}
and recursively define the symmetric product
\begin{equation}\label{eq:suzuki}
 S_{2k}(t)=S_{2k-2}(p_kt)^2
          S_{2k-2}((1-4p_k)t)
          S_{2k-2}(p_kt)^2.
\end{equation}
Then $\log S_{2k}(t)=t(X+Y)+O(t^{2k+1})$. For each fixed $k$,
$S_{2k}(t/n)^n$ has error $O(n^{-2k})$ in norm, locally uniformly
in $t$. Every term in its local logarithm has an explicit finite
multifactor word-cut expansion.
\end{theorem}
\begin{proof}
Inductively suppose that $S_{2k-2}$ is symmetric and its logarithm
is $tS+t^{2k-1}E+O(t^{2k+1})$, $S=X+Y$. The palindromic product
\eqref{eq:suzuki} is symmetric as well. Its linear terms add to
$tS$ because $4p_k+(1-4p_k)=1$. Through degree $2k-1$, all
commutators between the linear terms vanish, since they are scalar
multiples of $S$, and commutators involving $E$ have degree at
least $2k$. Thus the degree-$2k-1$ error is
\[
 \bigl(4p_k^{2k-1}+(1-4p_k)^{2k-1}\bigr)E=0,
\]
since $1-4p_k=-4^{1/(2k-1)}p_k$. The degree-$2k$ term also
vanishes by symmetry, leaving an error of order $2k+1$.
This establishes the induction from \eqref{eq:symmetriclog}.
For fixed $k$ the expression is a finite product of exponentials;
its logarithm therefore converges near zero. Scaling by $t/n$ and
raising to the $n$th power yields a logarithmic error
$O(n^{-2k})$, and the exponential estimate used in
\cref{thm:trotter} proves the norm assertion. Applying the
multifactor coefficient formula to this finite product gives the
all-orders expansion claimed.
\end{proof}
The middle coefficient $1-4p_k$ is negative. Such negative steps
are harmless for bounded operators but can be a serious obstruction
for unbounded semigroup generators. The hypothesis in this section
is therefore part of the result, not an optional technicality.

\section{Maurer--Cartan coordinates, vielbeins, and all-order metrics}
\label{sec:geometry}
\subsection{The infinitesimal formula and its full component expansion}
Let $e_1,\ldots,e_d$ be a basis of a finite-dimensional Lie algebra,
with structure constants
\begin{equation}\label{eq:structure}
 [e_i,e_j]=f_{ij}^{\ k}e_k,\qquad X=X^i e_i.
\end{equation}
Repeated upper--lower indices are summed. For a differentiable
Lie-algebra-valued function $X$, the differential formula already
proved gives the complete infinitesimal identity
\begin{equation}\label{eq:infinitesimal}
 e^{-X}de^X=\sum_{n\geq0}\frac{(-1)^n}{(n+1)!}\ad_X^n(dX)
 =dX-\tfrac12[X,dX]+\tfrac16[X,[X,dX]]-\cdots.
\end{equation}
In a nonmatrix Lie group, $e^{-X}de^X$ is the conventional notation
for the pullback of the left Maurer--Cartan form, so the intrinsic
proof in \cref{sec:localgroups} applies unchanged.
The displayed component terms on the reference page are
\begin{equation}\label{eq:infinitesimalcomponents}
 \begin{aligned}
 e^{-X}de^X={}&e_k\,dX^k
 -\tfrac12 X^i dX^j f_{ij}^{\ k}e_k
 +\tfrac16X^iX^j dX^k f_{jk}^{\ \ell}f_{i\ell}^{\ m}e_m\\
 &+\text{terms with at least three factors }X^i.
 \end{aligned}
\end{equation}
Each term follows by inserting \eqref{eq:structure} in the
successive commutators; thus its indices and signs are fixed by the
bracket convention rather than guessed from tensor notation.

Define the adjoint matrix and the left exponential-coordinate
coframe matrix by
\begin{equation}\label{eq:MW}
 M^a_{\ b}=X^i f_{ib}^{\ a},\qquad
 W(X)=\phi_-(M)=\sum_{n\geq0}\frac{(-1)^n}{(n+1)!}M^n.
\end{equation}
Then every component at every order is specified by
\begin{equation}\label{eq:thetaW}
 e^{-X}de^X=e_a\,W^a_{\ b}(X)dX^b,
\end{equation}
where for $n\geq1$
\begin{equation}\label{eq:componentall}
 (M^n)^a_{\ b}
 =\sum_{\substack{i_1,\ldots,i_n\\c_1,\ldots,c_{n-1}}}
 X^{i_1}\cdots X^{i_n}
 \prod_{r=1}^n f_{i_r c_r}^{\ c_{r-1}},
 \qquad c_0=a,\quad c_n=b.
\end{equation}
For $n=1$ this formula is simply $X^if_{ib}^{\ a}$, with the
intermediate-index string absent. Formula \eqref{eq:componentall}
is ordinary matrix multiplication written out, and the factorial
series in \eqref{eq:MW} converges for every finite matrix $M$.

\begin{proposition}[Which matrix is the vielbein?]\label{prop:vielbein}
The coframe coefficients of the Maurer--Cartan form in exponential
coordinates are $W^a_{\ b}$, not $M^a_{\ b}$. When the Lie algebra has positive dimension, $M$ is
always singular, whereas $W(0)=I$ and $W$ is invertible near zero.
More precisely,
\begin{equation}\label{eq:Wcriterion}
 W(X)\text{ is invertible}
 \quad\Longleftrightarrow\quad
 \spec(\ad_X)\cap(2\pi i\Z\setminus\{0\})=\varnothing,
\end{equation}
where the spectrum is taken after complexification for a real Lie
algebra. On this set $W^{-1}=f(M)$, interpreted by analytic
functional calculus, not necessarily by a convergent Taylor series
about zero at every such $M$.
\end{proposition}
\begin{proof}
Equation \eqref{eq:thetaW} identifies the coframe. If $X\ne0$, then
$MX=[X,X]=0$; at $X=0$, $M=0$. Thus $M$ cannot be an invertible
frame matrix. By contrast the constant term of $W$ is $I$.
Triangularize $M$ over $\C$. A power series in a triangular matrix
has diagonal entries equal to that scalar function at the diagonal
entries of $M$. Therefore
$\det W=\prod_{\lambda\in\spec M}\phi_-(\lambda)$,
with eigenvalues counted with multiplicity. The zeros of
$\phi_-$ are exactly the nonzero $2\pi i$ integers, proving
\eqref{eq:Wcriterion}. Where no zero is present, the analytic
identity $\phi_-f=1$ yields $W^{-1}=f(M)$.
\end{proof}
Some authors reserve ``vielbein'' for the dual vector frame rather
than its coframe. With that convention the frame coefficients are
$W^{-1}$, still not $M$. The quotient
$(I-e^{-M})M^{-1}$ used on the reference page is a shorthand for
the entire function $\phi_-(M)$; a literal inverse of $M$ does not
exist in this setting. This is a correction of a mathematical
misidentification on the page, not an alternative sign convention.

\subsection{Pullbacks and the hypotheses needed for a metric}
Choose a symmetric bilinear form $B$ on the Lie algebra, with
matrix $B_{ab}=B(e_a,e_b)$. Its left-invariant extension to the
group evaluates tangent vectors by first applying the
Maurer--Cartan form. For a smooth map $F:N\to G$, near any fixed $q_0$ a fixed
left translation replaces $F(q)$ by $F(q_0)^{-1}F(q)$ without
changing the pullback of a left-invariant tensor. We may therefore
work in an exponential neighborhood and write $F(q)=\exp X(q)$.
Set
\begin{equation}\label{eq:generalE}
 E^a_{\ \mu}(q)=W^a_{\ b}(X(q))\,\partial_\mu X^b(q).
\end{equation}
The complete pullback formula is then
\begin{equation}\label{eq:pullback}
 (F^*B)_{\mu\nu}=B_{ab}E^a_{\ \mu}E^b_{\ \nu}.
\end{equation}
This follows immediately from the definition of pullback: evaluate
$B$ on $F_*\partial_\mu,F_*\partial_\nu$ and use
\eqref{eq:thetaW}. In exponential coordinates themselves,
$X^b=q^b$, so the Jacobian is the identity and the page's formula
becomes
\begin{equation}\label{eq:metricW}
 g_{ij}=B_{ab}W^a_{\ i}W^b_{\ j},\qquad g=W^TBW.
\end{equation}
For an arbitrary map, the Jacobian factors in
\eqref{eq:generalE} cannot be omitted.

Several qualifications matter. A symmetric bilinear tensor need
not be a metric: a pseudo-Riemannian metric must be nondegenerate,
and a Riemannian metric must be positive definite. If $B$ is
positive definite and $F$ is an immersion, \eqref{eq:pullback} is
positive definite. If $B$ is indefinite, even an immersion can
have degenerate pullback: the immersion $q\mapsto(q,q)$ into
$\R^2$ with bilinear form $dx^2-dy^2$ pulls it back to zero. A
constant map pulls back every bilinear form to zero.

The Killing form is
\begin{equation}\label{eq:killing}
 B_K(u,v)=\tr(\ad_u\ad_v),\qquad
 (B_K)_{ab}=f_{ac}^{\ d}f_{bd}^{\ c}.
\end{equation}
It is symmetric by cyclicity of trace. Its invariance is proved by
\[
 B_K([x,u],v)+B_K(u,[x,v])
 =\tr\bigl([\ad_x,\ad_u]\ad_v+
             \ad_u[\ad_x,\ad_v]\bigr)=0,
\]
using $\ad_{[x,u]}=[\ad_x,\ad_u]$, which is Jacobi's identity.
But the Killing form is not nondegenerate in every Lie algebra:
in an abelian algebra every adjoint map is zero, so $B_K=0$.
Consequently calling its pullback a ``metric'' without additional
nondegeneracy hypotheses is not a valid general assertion.
The formula \eqref{eq:metricW} is valid regardless of whether it
defines a metric or only a possibly degenerate tensor.

\subsection{All-orders metric coefficients and their even resummation}
Without assuming invariance, substitution of \eqref{eq:MW} into
\eqref{eq:metricW} gives the absolutely convergent expansion
\begin{equation}\label{eq:metricdouble}
 g_{ij}=\sum_{p,q\geq0}\frac{(-1)^{p+q}}{(p+1)!(q+1)!}
 B_{ab}(M^p)^a_{\ i}(M^q)^b_{\ j}.
\end{equation}
If $B$ is invariant, the answer becomes considerably simpler.

\begin{theorem}[Complete invariant-metric expansion]\label{thm:metric}
For an invariant symmetric bilinear form $B$,
\begin{align}
 g&=B\frac{e^M+e^{-M}-2I}{M^2}
   =B\left(\frac{\sinh(M/2)}{M/2}\right)^2
       \label{eq:metricclosed}\\
  &=B\sum_{n\geq0}\frac{2M^{2n}}{(2n+2)!}
   =B\left(I+\frac{M^2}{12}+\frac{M^4}{360}
             +\frac{M^6}{20160}+\cdots\right).
       \label{eq:metriceven}
\end{align}
All quotients in this display denote entire functions with their
removable values at zero. No inverse of $M$ or $M^2$ is assumed.
\end{theorem}
\begin{proof}
Invariance says $B([X,u],v)+B(u,[X,v])=0$, or
$M^TB+BM=0$. Hence $(M^p)^TB=(-1)^pBM^p$, by induction.
The series for $W$ therefore gives
$W^TB=B\phi_-(-M)$. Multiply on the right by $W=\phi_-(M)$:
\[
 W^TBW=B\phi_-(-M)\phi_-(M)
 =B\frac{(e^M-I)(I-e^{-M})}{M^2}.
\]
The numerator is $e^M+e^{-M}-2I$. The scalar hyperbolic identity
and its power series prove both displays. All operations are
between convergent entire functions of the same matrix, so their
factors commute as used in the simplification. The resulting
matrix is symmetric because it was obtained as $W^TBW$.
\end{proof}
If $B$ and $W$ are invertible, the inverse tensor is explicitly
$g^{-1}=f(M)B^{-1}f(M)^T$. In a neighborhood of zero, expanding
each $f(M)$ using \eqref{eq:bernfinite} supplies all coefficients
of this inverse as well.

\subsection{A concrete three-dimensional resummation}
Identify the rotation Lie algebra with $\R^3$ and bracket
$[x,v]=x\times v$. Here $M v=x\times v$ and $r=|x|$.
The vector triple-product identity gives
$M^2v=x(x\cdot v)-r^2v$ and hence $M^3=-r^2M$.
Summing the exponential series by reducing powers with this
identity gives Rodrigues' formula
\[
 e^{-sM}=I-\frac{\sin(sr)}rM+
              \frac{1-\cos(sr)}{r^2}M^2.
\]
Integrating from $s=0$ to $1$ proves the complete coframe formula
\begin{equation}\label{eq:so3W}
 W=I-\frac{1-\cos r}{r^2}M+
          \frac{r-\sin r}{r^3}M^2,
\end{equation}
with its Taylor values at $r=0$.
The ordinary dot product is invariant because
$(x\times u)\cdot v+u\cdot(x\times v)=0$. For $r>0$, let
$P_\parallel=xx^T/r^2$ and $P_\perp=I-P_\parallel$.
On the parallel line $M=0$; on the perpendicular plane
$M^2=-r^2I$. Formula \eqref{eq:metricclosed} therefore gives
\begin{equation}\label{eq:so3metric}
 W^TW=P_\parallel+
       \left(\frac{\sin(r/2)}{r/2}\right)^2P_\perp,
 \qquad
 \det W=\left(\frac{\sin(r/2)}{r/2}\right)^2.
\end{equation}
The determinant follows either from the eigenvalues $0,ir,-ir$
of $M$ or directly from \eqref{eq:so3W}. Thus the coordinate
coframe is singular at nonzero radii $2\pi,4\pi,\ldots$, precisely
as \eqref{eq:Wcriterion} predicts. At zero the whole expression
has the unambiguous limit $I$, despite the individual projector
notation there being unnecessary.

\section{Quantum-mechanical identities with rigorous interpretations}
\label{sec:quantum}
The algebraic calculations in the reference page are correct in the
Heisenberg algebra. For operators on a Hilbert space, however,
unboundedness and domains cannot be suppressed. We first prove the
exponentiated identities in a concrete representation, then exhibit
why a commutator identity on a dense domain alone is insufficient.
This avoids using the Stone--von Neumann classification theorem,
which is mentioned but not stated on the reference page.

\subsection{Position, momentum, and the Weyl relations}
On $L^2(\R)$ fix $\hbar>0$. On the Schwartz space $\mathcal S(\R)$
let
\begin{equation}\label{eq:QP}
 (Qf)(x)=xf(x),\qquad (Pf)(x)=-i\hbar f'(x).
\end{equation}
Direct differentiation gives
\begin{equation}\label{eq:CCR}
 [Q,P]f=i\hbar f\qquad(f\in\mathcal S(\R)).
\end{equation}
The self-adjoint realizations are the usual multiplication operator
and the Fourier-transform conjugate multiplication operator.
Schwartz space is a core for these operators: cutoff and smoothing
approximate vectors in the graph norm of a multiplication operator,
and the Fourier transform preserves Schwartz space.
Their unitary exponentials act explicitly by
\begin{equation}\label{eq:UV}
 U(a)f(x)=e^{iax}f(x),\qquad
 V(b)f(x)=f(x+\hbar b),\qquad a,b\in\R.
\end{equation}
These assertions also follow from the multiplication spectral
calculus and the Fourier transform, respectively; neither requires
BCH. The functions \eqref{eq:UV} are unitary by a change of
variables and strongly continuous by density of smooth compactly
supported functions and dominated convergence.

\begin{theorem}[Weyl and symmetric exponential identities]
\label{thm:weyl}
In this representation,
\begin{align}
 e^{iaQ}e^{ibP}
   &=e^{\,i(aQ+bP-ab\hbar I/2)},\label{eq:quantumbch}\\
 e^{i(aQ+bP)}
   &=e^{iaQ/2}e^{ibP}e^{iaQ/2},\label{eq:quantumsymmetric}\\
 U(a)V(b)&=e^{-iab\hbar}V(b)U(a).\label{eq:weyl}
\end{align}
More generally, with $\mathcal W(a,b)=e^{i(aQ+bP)}$,
\begin{equation}\label{eq:weylgroup}
 \mathcal W(a,b)\mathcal W(c,d)
 =e^{i\hbar(bc-ad)/2}\mathcal W(a+c,b+d).
\end{equation}
\end{theorem}
\begin{proof}
Define the unitary operator
\begin{equation}\label{eq:Wexplicit}
 \mathcal W(a,b)f(x)=e^{ia(x+\hbar b/2)}f(x+\hbar b).
\end{equation}
For fixed $a,b$, the operators $\mathcal W(ta,tb)$ form a
strongly continuous one-parameter group: direct multiplication of
the phases proves the group law on this line. Its derivative at
$t=0$ on Schwartz functions is $i(aQ+bP)$. To identify its
self-adjoint generator without ambiguity, for $b\ne0$ multiplication
by $e^{iax^2/(2b\hbar)}$, denoted $T$, preserves Schwartz space
and satisfies
\[
 T^{-1}(bP)T=aQ+bP.
\]
Thus this differential operator has the self-adjoint realization
$T^{-1}(bP)T$, and direct conjugation of the translations gives
\eqref{eq:Wexplicit}. For $b=0$ it is the multiplication operator
$aQ$. This proves the claimed identification with the exponential.

Now $U(a)V(b)f=e^{iax}f(x+\hbar b)$, which is
$e^{-iab\hbar/2}\mathcal W(a,b)f$. Multiplying
$U(a/2)V(b)U(a/2)$ gives precisely \eqref{eq:Wexplicit}.
Comparing $U(a)V(b)$ with $V(b)U(a)$ proves \eqref{eq:weyl}.
Finally multiply two instances of \eqref{eq:Wexplicit}; after
subtracting the phase of $\mathcal W(a+c,b+d)$, the remaining
constant phase is $\hbar(bc-ad)/2$. This proves
\eqref{eq:weylgroup} on all of $L^2(\R)$.
\end{proof}
The formal BCH calculation agrees: for $A=iaQ$ and $B=ibP$,
$[A,B]=-iab\hbar I$ is central, so only its one-half term remains.
The preceding direct proof establishes that this formal calculation
is legitimate for the chosen self-adjoint realizations.

For any finite list $(a_j,b_j)$, induction in
\eqref{eq:weylgroup} yields the complete finite-factor formula
\begin{equation}\label{eq:weylmulti}
 \prod_{j=1}^m\mathcal W(a_j,b_j)
 =\exp\left(\frac{i\hbar}{2}
             \sum_{j<k}(b_ja_k-a_jb_k)\right)
   \mathcal W\left(\sum_ja_j,\sum_jb_j\right).
\end{equation}
This is \eqref{eq:centralmulti} in its unitary Heisenberg
representation, with no omitted higher-order terms.

\subsection{Why the infinitesimal commutator is not enough}
\begin{proposition}\label{prop:CCRcounterexample}
A relation $[Q,P]=iI$ on a dense common invariant domain, even for
self-adjoint $Q,P$, does not by itself imply the Weyl relations for
all real parameters.
\end{proposition}
\begin{proof}
Work on $L^2(0,1)$. Let $Q$ be multiplication by $x$, a bounded
self-adjoint operator, and let $P=-i\,d/dx$ with periodic
self-adjoint domain
$\{f\in H^1(0,1):f(0)=f(1)\}$. Its Fourier basis
$e^{2\pi inx}$ has eigenvalues $2\pi n$; equivalently it is the
self-adjoint diagonal operator with these real eigenvalues.
The dense subspace $C_c^\infty(0,1)$ is invariant under both
operators and has $[Q,P]=iI$ by differentiation. But
$e^{iP}=I$, since $e^{2\pi in}=1$ on every Fourier basis vector.
The Weyl relation with parameter $b=1$ would imply
$e^{iaQ}=e^{-ia}e^{iaQ}$ for every real $a$, which is false unless
$a\in2\pi\Z$. Thus the dense-domain commutator is insufficient.
The subspace used here is not a core for the periodic derivative;
that distinction is exactly one of the domain issues hidden by
purely formal notation.
\end{proof}
There cannot even be bounded operators $A,B$ with $[A,B]=cI$ for
$c\ne0$. Indeed induction gives $[A,B^n]=ncB^{n-1}$, so
\[
 n|c|\norm{B^{n-1}}\leq2\norm{A}\norm{B^n}.
\]
Iterating implies
$\norm{B^n}\geq n!|c|^n\norm{I}/(2\norm{A})^n$, contrary to
$\norm{B^n}\leq\norm{B}^n$ as $n\to\infty$.
This explains why the bounded-operator convergence theorems cannot
simply be applied to the canonical position and momentum pair.

\subsection{Creation and annihilation in Bargmann--Fock space}
Let $\mathcal F$ be the Hilbert space of entire functions with norm
\begin{equation}\label{eq:focknorm}
 \norm{F}_{\mathcal F}^2=
 \frac1\pi\int_{\C}|F(z)|^2e^{-|z|^2}\,d^2z.
\end{equation}
Polar integration gives
$\langle z^m,z^n\rangle=\delta_{mn}n!$: the angular integral is
zero for $m\ne n$, and the radial integral for $m=n$ is $n!$.
Consequently
\begin{equation}\label{eq:fockbasis}
 |n\rangle=\frac{z^n}{\sqrt{n!}},\qquad n\geq0,
\end{equation}
is an orthonormal basis. Completeness follows by expanding an
entire function in its Taylor series, integrating first on circles,
and then using monotone convergence for the sum of squared
coefficient norms. In particular polynomials are dense.
On this dense polynomial domain define
\begin{equation}\label{eq:creation}
 a=\partial_z,\qquad a^\dagger=z,
 \qquad[a,a^\dagger]=I,\quad[a^\dagger,a]=-I.
\end{equation}
We take the Hilbert inner product to be linear in its second argument.
The basis relations
$a|n\rangle=\sqrt n|n-1\rangle$ and
$a^\dagger|n\rangle=\sqrt{n+1}|n+1\rangle$ prove that the two
operators are adjoint there in the usual densely defined sense.

\begin{theorem}[Displacement and normal ordering]\label{thm:displacement}
For $v\in\C$, the displacement operator has the explicit unitary
action
\begin{equation}\label{eq:Dexplicit}
 (D(v)F)(z)=e^{-|v|^2/2}e^{vz}F(z-\bar v).
\end{equation}
It is the unitary exponential generated by
$v a^\dagger-\bar v a$, and on polynomials it satisfies the normal
ordering identity
\begin{equation}\label{eq:Dnormal}
 D(v)=e^{v a^\dagger-\bar v a}
     =e^{v a^\dagger}e^{-\bar v a}e^{-|v|^2/2}.
\end{equation}
The right-hand side extends by continuity from that domain to the
bounded operator $D(v)$; it is not a claim that both unbounded
factors can be composed on every vector without a domain choice.
For all $u,v\in\C$,
\begin{equation}\label{eq:Dproduct}
 D(v)D(u)=
 \exp\left(\frac{v\bar u-u\bar v}{2}\right)D(v+u).
\end{equation}
\end{theorem}
\begin{proof}
In \eqref{eq:focknorm}, substitute $z=w+\bar v$ in the norm of
\eqref{eq:Dexplicit}. The total real exponent simplifies to
\[
 -|v|^2+2\operatorname{Re}(v(w+\bar v))-|w+\bar v|^2=-|w|^2.
\]
Thus $D(v)$ is isometric. Direct substitution gives $D(-v)$ as its
inverse, so it is unitary. Multiplying its explicit actions yields
\[
 D(v)D(u)F(z)=e^{-(|v|^2+|u|^2)/2-u\bar v}
            e^{(v+u)z}F(z-\bar v-\bar u).
\]
Comparison with $D(v+u)$ proves \eqref{eq:Dproduct}, including its
sign and factor $1/2$.

For a fixed $v$, $D(tv)$ is a strongly continuous one-parameter
unitary group for real $t$; its group law follows from
\eqref{eq:Dproduct} because the central phase vanishes on this
line. Strong continuity follows first for polynomials from the
explicit formula and Gaussian integrability, and then for every
vector by density and unitarity. Its derivative at zero on
polynomials is $v a^\dagger-\bar v a$. The exponential notation in
\eqref{eq:Dnormal} denotes this skew-adjoint generator.
There is no ambiguity about its closure: on a basis vector the
$k$th power $T^k$, $T=v a^\dagger-\bar v a$, satisfies
\begin{equation}\label{eq:analyticvectors}
 \norm{T^k|n\rangle}\leq(2|v|)^k
              \sqrt{\frac{(n+k)!}{n!}}.
\end{equation}
This follows by noting that at step $j$ all occupied levels are at
most $n+j$, and each of $a,a^\dagger$ has norm at most the square
root of the next level on that finite-dimensional subspace.
Thus $\sum_k\norm{T^kF}|t|^k/k!$ converges for every polynomial
$F$ and every $t$.
To check uniqueness explicitly, if a vector $\eta$ were in a
nonzero deficiency space for the skew-symmetric operator $T$, it
would satisfy $T^*\eta=\eta$ or $T^*\eta=-\eta$.
For polynomials the identity
$D(tv)F=\sum_{k\geq0}t^kT^kF/k!$ follows from the integral Taylor
remainder: its norm is at most
$|t|^{k+1}\norm{T^{k+1}F}/(k+1)!$, which tends to zero by
\eqref{eq:analyticvectors}. Hence
$\langle\eta,D(tv)F\rangle$ would equal
$e^t\langle\eta,F\rangle$ or
$e^{-t}\langle\eta,F\rangle$. Unitarity bounds this scalar
for all positive and negative real $t$, forcing
$\langle\eta,F\rangle=0$ for all polynomials and hence $\eta=0$.
To finish the operator argument explicitly, let $S$ be the closure
of $T$, which exists because $T$ is densely defined and
skew-symmetric. The identity
$\norm{(I\pm S)u}^2=\norm{u}^2+\norm{Su}^2$ shows that both
ranges are closed. Their orthogonal complements are the two
vanishing deficiency kernels, so both ranges are the entire
Hilbert space. For $\eta\in\operatorname{Dom}(S^*)$, solve
$(I-S)u=(I+S^*)\eta$. Since $S^*u=-Su$, the difference $\eta-u$
lies in $\ker(I+S^*)=0$. Thus $\eta=u\in\operatorname{Dom}(S)$,
proving $S^*=-S$. The explicit unitary group therefore supplies
this unique skew-adjoint closure, with polynomials as a core.

Finally, on a polynomial $F$, $e^{-\bar v a}F=F(z-\bar v)$ by the
finite Taylor formula, while $e^{v a^\dagger}$ multiplies by
$e^{vz}$. Their ordered product with the scalar factor is exactly
\eqref{eq:Dexplicit}. This proves the normal ordering assertion
on the stated domain and hence its bounded extension.
\end{proof}
The commutator of the two generators is
\[
 [v a^\dagger-\bar v a,\ u a^\dagger-\bar u a]
 =(v\bar u-u\bar v)I.
\]
It is central and purely imaginary. Thus \eqref{eq:Dproduct} is
also exactly the central-commutator BCH identity, now justified by
a representation-level proof rather than by manipulating undefined
operator products.

\subsection{All matrix elements and finite products}
Every finite displacement product has the closed expression
\begin{equation}\label{eq:Dmulti}
 \prod_{j=1}^mD(v_j)
 =\exp\left(\frac12\sum_{j<k}
           (v_j\bar v_k-v_k\bar v_j)\right)
        D\left(\sum_jv_j\right),
\end{equation}
by induction in \eqref{eq:Dproduct}. Acting on the vacuum gives the
full coherent-state series
\begin{equation}\label{eq:coherent}
 |v\rangle=D(v)|0\rangle
 =e^{-|v|^2/2}\sum_{n\geq0}\frac{v^n}{\sqrt{n!}}|n\rangle,
 \qquad
 \langle u|v\rangle=
 e^{-(|u|^2+|v|^2)/2+\bar u v}.
\end{equation}
The first formula is the exponential series for $e^{vz}$, and the
second follows by the orthonormal basis expansion and the scalar
exponential series. The norm at $u=v$ is one.
More generally all number-state matrix elements are the finite sums
\begin{equation}\label{eq:Dmatrix}
 \boxed{\displaystyle
 \langle m|D(v)|n\rangle
 =e^{-|v|^2/2}\sqrt{m!n!}
   \sum_{k=0}^{\min(m,n)}
   \frac{v^{m-k}(-\bar v)^{n-k}}
        {(m-k)!(n-k)!k!}.}
\end{equation}
To derive this, apply the finite lowering series
$e^{-\bar v a}$ to $|n\rangle$, leaving level $k$ with coefficient
$(-\bar v)^{n-k}\sqrt{n!/k!}/(n-k)!$.
The raising series $e^{v a^\dagger}$ then reaches level $m$ with
coefficient $v^{m-k}\sqrt{m!/k!}/(m-k)!$. Multiplying and summing
over the admissible $k$ proves \eqref{eq:Dmatrix}.
These formulas provide all terms of the oscillator applications;
no asymptotic truncation remains.

\appendix
\section{An elementary PBW proof and the enveloping Hopf algebra}
\label{sec:pbw}
This appendix supplies the injectivity statement used in the
universal-algebra discussion. It also proves that the Hopf structure
exists on an arbitrary universal enveloping algebra, not only on
the free associative algebra.

Let $\h$ be a Lie algebra with an ordered basis $\{e_i\}_{i\in I}$.
Its enveloping algebra is the quotient of the tensor algebra on
its underlying vector space by the relations
\begin{equation}\label{eq:Urelation}
 uv-vu=[u,v]\qquad(u,v\in\h).
\end{equation}
It has the following universal property: every Lie homomorphism
from $\h$ into the commutator algebra of an associative algebra
extends uniquely to an associative homomorphism from $U(\h)$.
Indeed it first extends uniquely to the tensor algebra, and
\eqref{eq:Urelation} is exactly the condition that the resulting
map annihilate the quotient ideal.

\begin{theorem}[Poincar\'e--Birkhoff--Witt]\label{thm:pbw}
The ordered monomials
$e_{i_1}\cdots e_{i_m}$ with $i_1\leq\cdots\leq i_m$, including
the empty monomial, form a vector-space basis of $U(\h)$.
In particular the natural map $\h\to U(\h)$ is injective.
\end{theorem}
\begin{proof}
Replace an adjacent inverted pair by
\begin{equation}\label{eq:rewrite}
 e_je_i\longmapsto e_ie_j+[e_j,e_i]\qquad(j>i),
\end{equation}
expanding the bracket in the given basis. Order words by their
length first and their number of inversions second. The swapped
word has the same length and one fewer inversion, whereas every
bracket term has smaller length. This is a well-founded decreasing
measure, so repeated rewriting terminates in a linear combination
of ordered words. It proves spanning.

To prove uniqueness of the result, consider two first reduction
choices. Reductions on disjoint adjacent pairs commute. The only
overlapping ambiguity is a word $cba$ with three strictly decreasing
basis indices $c>b>a$. Reducing the left inverted pair first and
then ordering the remaining triple gives
\[
 abc+[b,a]c+b[c,a]+[c,b]a.
\]
Reducing the right inverted pair first instead gives
\[
 abc+a[c,b]+[c,a]b+c[b,a].
\]
The difference, using the length-two relations, reduces to
\[
 [[b,a],c]+[b,[c,a]]+[[c,b],a]=0
\]
by Jacobi. Each bracket-expanded term has smaller length than the
original triple. With a prefix and suffix attached, it still has
smaller length than the original word. Induction on the decreasing
measure therefore already guarantees a unique normal form for
these differences.

More explicitly, suppose every strictly smaller word has a unique
normal form. For two reductions of the current word, the disjoint
case has a common next reduction, while the overlapping case just
computed has difference with zero normal form among smaller words.
All subsequent choices act on smaller words, so their normal forms
agree by induction. This proves uniqueness for the current word,
and hence for every word and by linearity for every polynomial.
The resulting normal-form map annihilates each contextual relation
\eqref{eq:rewrite}, hence its two-sided ideal. Conversely a
polynomial and its normal form differ by that ideal. Therefore the
quotient is exactly the vector space freely spanned by the ordered
words. In particular the length-one basis vectors remain linearly
independent, proving injectivity.
\end{proof}
No matrix realization of $\h$ and no BCH theorem was used in this
proof, so the earlier use of PBW is not circular.

On the tensor algebra set
$\Delta u=u\otimes1+1\otimes u$, $\eps(u)=0$, and $S(u)=-u$
for $u\in\h$, extending $\Delta,\eps$ multiplicatively and $S$
antimultiplicatively. For a defining relation
$r=uv-vu-[u,v]$, direct expansion gives
\[
 \Delta r=r\otimes1+1\otimes r,\qquad \eps(r)=0,\qquad S(r)=-r.
\]
Consequently these operations descend to the quotient $U(\h)$.
The coproduct is coassociative on generators and therefore on the
whole algebra; the counit and antipode identities proved for tensor
algebras descend as well. This is the natural enveloping Hopf
algebra structure asserted on the reference page.

To see concretely why an infinite-series coproduct needs a
completed tensor product, restrict to one generator and consider
$F(X)=\sum_{n\geq0}X^n$. Its formal coproduct has coefficient matrix
\[
 [X^p\otimes X^q]\Delta F=\binom{p+q}{p}\qquad(p,q\geq0).
\]
The leading $(N+1)\times(N+1)$ matrix is $PP^T$, where
$P_{pk}=\binom pk$ for $k\leq p$ and is zero otherwise. The identity
$\binom{p+q}{p}=\sum_k\binom pk\binom qk$ is Vandermonde's
identity, obtained by comparing coefficients in
$(1+t)^p(1+t)^q$. Since $P$ is triangular with diagonal $1$,
every such matrix has determinant $1$. Its rank is therefore
unbounded with $N$. But the coefficient matrix of a finite sum of
separated tensors $\sum_{j=1}^m A_j(X)\otimes B_j(X)$ has rank
at most $m$. Thus $\Delta F$ is not in the ordinary algebraic
tensor product, while it is in the completed tensor product used
in \cref{sec:hopf}.

\section{Complete word certificates for the displayed fifth and sixth degrees}
\label{sec:certificate}
In the following tables an entry is simultaneously the coefficient
$c(w)$ computed by the finite cut formula \eqref{eq:wordcuts} and
the coefficient obtained by expanding \eqref{eq:z5} or
\eqref{eq:z6}. Thus the tables specify every coordinate needed for
the polynomial identities, including zeros. Right-nested brackets
are expanded by the rule
$\RB{av}=a\RB{v}-\RB{v}a$, starting with $\RB{a}=a$.
For example, the coefficient of $X^4Y$ in $Z_5$ is $-1/720$,
and that of $X^5Y$ in $Z_6$ is zero, as the single-$Y$ Bernoulli
coefficients also predict.

\subsection*{Degree 5 word coefficients}
\begin{center}\small\begin{tabular}{lrlr}\toprule
Word & Coefficient & Word & Coefficient\\\midrule
\texttt{XXXXX} & $0$ & \texttt{YXXXX} & $-\tfrac{1}{720}$\\
\texttt{XXXXY} & $-\tfrac{1}{720}$ & \texttt{YXXXY} & $\tfrac{1}{180}$\\
\texttt{XXXYX} & $\tfrac{1}{180}$ & \texttt{YXXYX} & $-\tfrac{1}{120}$\\
\texttt{XXXYY} & $\tfrac{1}{180}$ & \texttt{YXXYY} & $-\tfrac{1}{120}$\\
\texttt{XXYXX} & $-\tfrac{1}{120}$ & \texttt{YXYXX} & $-\tfrac{1}{120}$\\
\texttt{XXYXY} & $-\tfrac{1}{120}$ & \texttt{YXYXY} & $\tfrac{1}{30}$\\
\texttt{XXYYX} & $-\tfrac{1}{120}$ & \texttt{YXYYX} & $-\tfrac{1}{120}$\\
\texttt{XXYYY} & $\tfrac{1}{180}$ & \texttt{YXYYY} & $\tfrac{1}{180}$\\
\texttt{XYXXX} & $\tfrac{1}{180}$ & \texttt{YYXXX} & $\tfrac{1}{180}$\\
\texttt{XYXXY} & $-\tfrac{1}{120}$ & \texttt{YYXXY} & $-\tfrac{1}{120}$\\
\texttt{XYXYX} & $\tfrac{1}{30}$ & \texttt{YYXYX} & $-\tfrac{1}{120}$\\
\texttt{XYXYY} & $-\tfrac{1}{120}$ & \texttt{YYXYY} & $-\tfrac{1}{120}$\\
\texttt{XYYXX} & $-\tfrac{1}{120}$ & \texttt{YYYXX} & $\tfrac{1}{180}$\\
\texttt{XYYXY} & $-\tfrac{1}{120}$ & \texttt{YYYXY} & $\tfrac{1}{180}$\\
\texttt{XYYYX} & $\tfrac{1}{180}$ & \texttt{YYYYX} & $-\tfrac{1}{720}$\\
\texttt{XYYYY} & $-\tfrac{1}{720}$ & \texttt{YYYYY} & $0$\\
\bottomrule\end{tabular}\end{center}
\subsection*{Degree 6 word coefficients}
\begin{center}\small\begin{tabular}{lrlrlrlr}\toprule
Word & Coefficient & Word & Coefficient & Word & Coefficient & Word & Coefficient\\\midrule
\texttt{XXXXXX} & $0$ & \texttt{XYXXXX} & $0$ & \texttt{YXXXXX} & $0$ & \texttt{YYXXXX} & $\tfrac{1}{1440}$\\
\texttt{XXXXXY} & $0$ & \texttt{XYXXXY} & $\tfrac{1}{360}$ & \texttt{YXXXXY} & $0$ & \texttt{YYXXXY} & $0$\\
\texttt{XXXXYX} & $0$ & \texttt{XYXXYX} & $0$ & \texttt{YXXXYX} & $-\tfrac{1}{360}$ & \texttt{YYXXYX} & $\tfrac{1}{240}$\\
\texttt{XXXXYY} & $-\tfrac{1}{1440}$ & \texttt{XYXXYY} & $-\tfrac{1}{240}$ & \texttt{YXXXYY} & $0$ & \texttt{YYXXYY} & $0$\\
\texttt{XXXYXX} & $0$ & \texttt{XYXYXX} & $0$ & \texttt{YXXYXX} & $\tfrac{1}{240}$ & \texttt{YYXYXX} & $\tfrac{1}{240}$\\
\texttt{XXXYXY} & $\tfrac{1}{360}$ & \texttt{XYXYXY} & $\tfrac{1}{60}$ & \texttt{YXXYXY} & $0$ & \texttt{YYXYXY} & $0$\\
\texttt{XXXYYX} & $0$ & \texttt{XYXYYX} & $0$ & \texttt{YXXYYX} & $\tfrac{1}{240}$ & \texttt{YYXYYX} & $\tfrac{1}{240}$\\
\texttt{XXXYYY} & $\tfrac{1}{360}$ & \texttt{XYXYYY} & $\tfrac{1}{360}$ & \texttt{YXXYYY} & $0$ & \texttt{YYXYYY} & $0$\\
\texttt{XXYXXX} & $0$ & \texttt{XYYXXX} & $0$ & \texttt{YXYXXX} & $-\tfrac{1}{360}$ & \texttt{YYYXXX} & $-\tfrac{1}{360}$\\
\texttt{XXYXXY} & $-\tfrac{1}{240}$ & \texttt{XYYXXY} & $-\tfrac{1}{240}$ & \texttt{YXYXXY} & $0$ & \texttt{YYYXXY} & $0$\\
\texttt{XXYXYX} & $0$ & \texttt{XYYXYX} & $0$ & \texttt{YXYXYX} & $-\tfrac{1}{60}$ & \texttt{YYYXYX} & $-\tfrac{1}{360}$\\
\texttt{XXYXYY} & $-\tfrac{1}{240}$ & \texttt{XYYXYY} & $-\tfrac{1}{240}$ & \texttt{YXYXYY} & $0$ & \texttt{YYYXYY} & $0$\\
\texttt{XXYYXX} & $0$ & \texttt{XYYYXX} & $0$ & \texttt{YXYYXX} & $\tfrac{1}{240}$ & \texttt{YYYYXX} & $\tfrac{1}{1440}$\\
\texttt{XXYYXY} & $-\tfrac{1}{240}$ & \texttt{XYYYXY} & $\tfrac{1}{360}$ & \texttt{YXYYXY} & $0$ & \texttt{YYYYXY} & $0$\\
\texttt{XXYYYX} & $0$ & \texttt{XYYYYX} & $0$ & \texttt{YXYYYX} & $-\tfrac{1}{360}$ & \texttt{YYYYYX} & $0$\\
\texttt{XXYYYY} & $-\tfrac{1}{1440}$ & \texttt{XYYYYY} & $0$ & \texttt{YXYYYY} & $0$ & \texttt{YYYYYY} & $0$\\
\bottomrule\end{tabular}\end{center}


The general formula is not inferred from these tables; it was
proved in \cref{thm:wordcuts,thm:dynkin}. The tables verify the
specific low-degree displays and make every arithmetic step
independently reproducible. The accompanying TSV and JSON files
continue the complete nonzero associative coefficients through
degree ten. Their size is modest because repeated word coefficients
have already been combined.

\section{Reproducible exact computations}
\label{sec:reproducibility}
The program \texttt{code/verify\_bch.py} uses Python's standard
library only, with exact rational numbers throughout. A polynomial
is a dictionary from words to rational coefficients, and products
concatenate words while discarding degrees above a chosen cutoff.
The empty word represents the identity. This works in the finite
quotient of the free associative algebra by words of degree greater
than the cutoff; no sampling of numerical matrices is involved.

The core logarithmic calculation is simply
\begin{lstlisting}[language=Python]
E = exp(X) * exp(Y)             # all operations truncated by degree
U = E - 1
Z = sum((-1)**(k-1) * U**k / k for k in range(1, N+1))
\end{lstlisting}
The actual implementation retains exact fractions and sparse word
dictionaries. It compares this result with the independent
Bernoulli differential recurrence, the right Dynkin projection,
and the independently enumerated cut formula. Three recursive
Zassenhaus constructions are compared with explicit enumeration
of the finite weighted-tree formula. Their products are then
multiplied out and compared with $e^{X+Y}$.

The included run used cutoff $N=10$ and passed all 20 named checks.
Most checks run through degree ten; two independent exhaustive
checks, the word-cut enumeration and unshuffle coproduct, run
through degree eight, and three-letter associativity through
degree six. The page's finite displays are checked through their
full displayed degrees, six for BCH and four for Zassenhaus.
These exact finite checks do not claim machine verification of
the infinite-order or analytic proofs.

\begin{center}
\begin{tabular}{rrrr}
\toprule
Degree & Nonzero BCH words & Nonzero $C_n$ words & Weighted trees\\
\midrule
1&2&0&0\\
2&2&2&1\\
3&6&6&1\\
4&4&12&2\\
5&30&30&2\\
6&28&42&5\\
7&126&126&6\\
8&124&234&16\\
9&390&510&22\\
10&388&954&54\\
\bottomrule
\end{tabular}
\end{center}
These counts concern the associative word expansion, not dimensions
of homogeneous free Lie spaces and not the number of terms in a
chosen Hall basis. The weighted-tree column counts trees in
\cref{thm:trees} before their associative terms are combined.

To reproduce the results from the archive's top-level directory:
\begin{lstlisting}[language=bash]
python3 code/verify_bch.py --degree 10 --output data
latexmk -pdf -interaction=nonstopmode -halt-on-error bch_complete.tex
\end{lstlisting}
Python 3.10 or later is suitable. The verifier permits degrees
between six and twelve; increasing the cutoff increases the amount
of exact arithmetic substantially. The LaTeX source uses standard
packages and does not require bibliography software or shell escape.
When \texttt{latexmk} is unavailable, run \texttt{pdflatex} twice
(or until cross-reference warnings disappear).

\section{Coverage index and necessary corrections}
\label{sec:coverage}
This index uses the pinned reference-page version stated in
\cref{sec:scope}. Every displayed mathematical formula and every
substantive mathematical assertion on that page is covered below.
A corrected assertion is identified as such rather than presented
as though an incorrect unrestricted version had been proved.

\small
\begin{longtable}{@{}p{0.47\textwidth}p{0.48\textwidth}@{}}
\toprule
Reference-page formula or assertion & Proof and interpretation here\\
\midrule
\endfirsthead
\toprule
Reference-page formula or assertion & Proof and interpretation here\\
\midrule
\endhead
\bottomrule
\endfoot
Product of exponentials has a formal Lie-series logarithm;
leading BCH terms &
\Cref{lem:explog,cor:formalBCH}; \eqref{eq:z1}--\eqref{eq:z3}.\\[4pt]
Only commutators and their nestings occur beyond $X+Y$ &
\Cref{cor:formalBCH}, including the ideal generated by $[X,Y]$.\\[4pt]
Small elements give a convergent BCH series &
\Cref{thm:convergence,cor:remainder}; the Banach Lie-algebra bound
also covers all finite-dimensional Lie algebras.\\[4pt]
Every group element near the identity is an exponential;
local multiplication is Lie-algebraic &
\Cref{thm:localgroup}; intrinsic differential and inverse function
theorem.\\[4pt]
Lie correspondence and integration of representations as applications &
\Cref{thm:localcorrespondence,thm:integration}; the latter includes
the simply connected hypothesis for global integration.\\[4pt]
Dynkin's multiple sum, its factorial and degree denominators &
\Cref{thm:dynkin}; every word coefficient in
\cref{thm:wordcuts}.\\[4pt]
Right-nesting notation and zero final blocks &
\eqref{eq:rightbracket}; the final assertion of
\cref{thm:dynkin}.\\[4pt]
Every BCH term through degree six &
\eqref{eq:z1}--\eqref{eq:z6}; derivation in
\cref{sec:lowdegree} and full certificates in
\cref{sec:certificate}.\\[4pt]
$X\leftrightarrow Y$ symmetry in alternating degrees &
\eqref{eq:swap}--\eqref{eq:swapdegree}, proved by inverse and
logarithm uniqueness.\\[4pt]
Poincar\'e integral formula &
\Cref{thm:poincare}; derived from the differential equation rather
than assumed as an external identity.\\[4pt]
$\psi(x)=x\log x/(x-1)$ and its complete Taylor series &
\eqref{eq:psiseries}; valid near $x=1$ or formally.\\[4pt]
$\psi(e^y)$ and normalized Bernoulli coefficients in the note &
\eqref{eq:psibernoulli}, \cref{prop:bernoulli}; convention
$B_1^+=1/2$ is explicit.\\[4pt]
Matrix tangent Lie algebra, commutator, and matrix exponential &
\Cref{sec:localgroups}; the flow equation identifies the ordinary
exponential.\\[4pt]
Associative logarithmic block sum &
\eqref{eq:associativeblocks}; absolute convergence in
\cref{thm:convergence}.\\[4pt]
Conditions involving $\log2$, including the note on
Hilbert--Schmidt norm &
\Cref{sec:convergence}; submultiplicativity is proved and the extra
bound on a constructed $Z$ is explained as unnecessary.\\[4pt]
Associative terms of degrees one through four &
\eqref{eq:assoclow}, obtained by direct bracket expansion.\\[4pt]
The trace logarithm identity &
\Cref{prop:trace}; local BCH branch required, with the exact
$2\pi i\Z$ ambiguity for other logarithms.\\[4pt]
The $\mathfrak{sl}_2(\C)$ example and its nonconvergence &
\Cref{thm:obstruction}; all matrix logarithms are explicitly
classified, and Abel's argument rules out even conditional BCH
convergence.\\[4pt]
Nonisomorphic Lie groups may have the same Lie algebra &
The explicit $\R$ and circle example after
\cref{thm:localgroup}.\\[4pt]
Sufficient norm bound $(\log2)/2$ &
\eqref{eq:dynkinbound}; it controls the unreduced nested multiple
sum.\\[4pt]
Commuting and central-commutator special cases; no smallness
restriction in the latter &
\Cref{thm:central}; a global differential proof in the stated
bounded/group settings.\\[4pt]
Expansion linear in $Y$, including the hyperbolic-cotangent form &
\eqref{eq:linearY}; all higher powers are supplied by
\cref{prop:Yall}.\\[4pt]
$[X,Y]=sY$ eliminates terms with two or more $Y$'s &
The Lie-monomial induction in \cref{thm:eigen}.\\[4pt]
Closed eigenvector BCH coefficient, resonance exclusion, and
$s=0$ limit &
\eqref{eq:eigenbch}; both removable zero and genuine nonzero
resonances are proved.\\[4pt]
Eigenvector braiding and adjoint dilation &
\eqref{eq:eigenbraid} and its proof; multiplying by $e^{-X}$
gives the displayed dilation.\\[4pt]
Formal exponential/logarithm bijection &
\Cref{lem:explog}.\\[4pt]
Coproduct of the generators and extension to series &
\eqref{eq:coproductgenerators}--\eqref{eq:unshuffle}; the tensor
product is completed. \Cref{sec:pbw} gives a counterexample to an
uncompleted interpretation.\\[4pt]
Exponential is grouplike iff exponent is primitive; grouplikes
form a group &
\Cref{lem:grouplike}, with constant term $1$ in the definition of
grouplike.\\[4pt]
Friedrichs' characterization and the primitive-logarithm existence
argument &
\Cref{lem:dsw,thm:friedrichs,cor:formalBCH}.\\[4pt]
Enveloping algebra of the free Lie algebra is the free associative
algebra; natural Hopf structure &
End of \cref{sec:hopf} and \cref{sec:pbw}; the completion is
specified explicitly.\\[4pt]
Zassenhaus product and every displayed exponent &
\Cref{thm:zassenhaus}; \eqref{eq:c2}--\eqref{eq:c4}.\\[4pt]
All higher Zassenhaus exponents are homogeneous Lie polynomials;
recursive determination &
\Cref{thm:zassenhaus}; residual, differential, and forward
algorithms; finite all-terms formula in \cref{thm:trees}.\\[4pt]
Alternative Zassenhaus notation starting at $n=1$ &
The paragraph after \cref{thm:zassenhaus} supplies the necessary
convention $C_1=Y$.\\[4pt]
Suzuki--Trotter consequence &
\Cref{thm:trotter,thm:suzuki}, including all-order local logarithmic
error series.\\[4pt]
$\Ad_{e^X}=e^{\ad_X}$, complete iterated-commutator sum, and its
linear differential equation &
\Cref{thm:hadamard} and \eqref{eq:adbinomial}.\\[4pt]
Conjugation of $e^Y$ and general braiding identity &
\eqref{eq:conjexp}--\eqref{eq:braiding}.\\[4pt]
The central-commutator application obtained by differentiating
$e^{sX}e^{sY}$ &
Proof of \cref{thm:central}, including its global initial-value
argument.\\[4pt]
Infinitesimal $e^{-X}de^X$ and its component terms &
\eqref{eq:infinitesimal}--\eqref{eq:infinitesimalcomponents};
all component orders in \eqref{eq:componentall}.\\[4pt]
$M^a_{\ b}=X^if_{ib}^{\ a}$, series for $W$, and quotient notation &
\eqref{eq:MW}; \cref{prop:vielbein} explains the removable
functional calculus and singularity of $M$.\\[4pt]
The page's statement that $M$ is a vielbein &
\emph{Corrected}: $W$ is the coframe matrix and $W^{-1}$ its dual
frame, where defined; $M$ is always singular.\\[4pt]
Pullback tensor $g_{ij}=W^a_{\ i}W^b_{\ j}B_{ab}$ &
\eqref{eq:pullback}--\eqref{eq:metricW}; an arbitrary map needs
its Jacobian. Full series in \cref{thm:metric}.\\[4pt]
Killing form and assertions that the pullback is a metric &
\emph{Qualified}: \eqref{eq:killing} and its invariance proof;
nondegeneracy and pullback hypotheses, with explicit counterexamples,
in \cref{sec:geometry}.\\[4pt]
Canonical position--momentum commutator &
\eqref{eq:QP}--\eqref{eq:CCR}, proved on Schwartz space.\\[4pt]
Exponentiated position--momentum and symmetric identities &
\Cref{thm:weyl}; direct unitary actions establish the domain-safe
meaning. \Cref{prop:CCRcounterexample} explains the limitation of a
bare infinitesimal relation.\\[4pt]
$[a^\dagger,a]=-I$ and displacement normal ordering &
\eqref{eq:creation} and \cref{thm:displacement}; explicitly stated
dense domain and bounded extension.\\[4pt]
Product of two displacements and its central phase &
\eqref{eq:Dproduct}; all finite products and number-state matrix
elements in \eqref{eq:Dmulti}--\eqref{eq:Dmatrix}.\\[4pt]
Heisenberg nilpotence underlying those applications &
\eqref{eq:heisenberglaw}, \cref{prop:nilpotent}, and the displayed
central generator commutator after \cref{thm:displacement}.\\
\end{longtable}
\normalsize

The genuinely all-terms outputs can be located without scanning
the index: every BCH coefficient is in \eqref{eq:wordcuts} and
\eqref{eq:allbidegrees}; every power of $Y$ in
\eqref{eq:Yrecurrence} and \eqref{eq:Yclosed}; every Zassenhaus
exponent in \eqref{eq:treeformula} and \eqref{eq:treelie}; every
Maurer--Cartan component in \eqref{eq:MW}--\eqref{eq:componentall};
every invariant metric coefficient in \eqref{eq:metriceven};
and every oscillator matrix element in \eqref{eq:Dmatrix}.
These statements specify coefficients at arbitrary order, not just
an extensive initial truncation.

\clearpage
\begin{thebibliography}{99}
\bibitem{wiki}
Wikipedia contributors,
\emph{Baker--Campbell--Hausdorff formula}, revision 1368909113,
11 August 2026, consulted 16 September 2026.
\url{https://en.wikipedia.org/w/index.php?oldid=1368909113}.
The page specifies the scope of the present paper, not a source
substituted for any proof.

\bibitem{dynkin}
E. B. Dynkin,
Calculation of the coefficients in the Campbell--Hausdorff formula
(in Russian),
\emph{Doklady Akademii Nauk SSSR} \textbf{57} (1947), 323--326.

\bibitem{casas}
F. Casas, A. Murua, and M. Nadinic,
Efficient computation of the Zassenhaus formula,
\emph{Computer Physics Communications} \textbf{183} (2012),
2386--2391.
\href{https://doi.org/10.1016/j.cpc.2012.06.006}{doi:10.1016/j.cpc.2012.06.006}.
\href{https://arxiv.org/abs/1204.0389}{arXiv:1204.0389}.

\bibitem{biagi}
S. Biagi, A. Bonfiglioli, and M. Matone,
On the Baker--Campbell--Hausdorff theorem: non-convergence and
prolongation issues,
\emph{Linear and Multilinear Algebra} \textbf{68} (2020),
1310--1328; published online 2018.
\href{https://doi.org/10.1080/03081087.2018.1540534}{doi:10.1080/03081087.2018.1540534}.
\href{https://arxiv.org/abs/1805.10089}{arXiv:1805.10089}.

\bibitem{trotter}
H. F. Trotter,
On the product of semi-groups of operators,
\emph{Proceedings of the American Mathematical Society}
\textbf{10} (1959), 545--551.
\href{https://doi.org/10.1090/S0002-9939-1959-0108732-6}{doi:10.1090/S0002-9939-1959-0108732-6}.

\bibitem{suzuki}
M. Suzuki,
General theory of fractal path integrals with applications to
many-body theories and statistical physics,
\emph{Journal of Mathematical Physics} \textbf{32} (1991),
400--407.
\href{https://doi.org/10.1063/1.529425}{doi:10.1063/1.529425}.
\end{thebibliography}
\end{document}
